%% 4장 행렬 해석 (새 목차, 2026-11-25)
\chapter{Matrix Analysis}
\label{ch:matrix-analysis}

\abstract{컴퓨터 안의 행렬은 처음부터 반올림된 값이므로 수치 선형대수의 첫 질문은 행렬이 조금 바뀌면 답이 얼마나 바뀌는가이다. 이 장은 노름, 특이값 분해, 조건수, 고유값과 특이값의 섭동, Schur 분해, Kronecker 곱, 양정치 행렬을 차례로 세운다. 구현으로는 넘치지 않는 노름, 선 합 노름, 수치적 랭크, Bartels--Stewart, \code{kron}·\code{kronsum}, Cholesky 판정을 본다.}


\section{Vector Norms}
\label{sec:vector-norms}

\subsection{Definition and Equivalence}

벡터공간 $\mathbb{F}^n$($\mathbb{F}$는 $\mathbb{R}$ 또는 $\mathbb{C}$) 위의 \term{노름}{norm}은 양의 정부호성($\|x\|=0\iff x=0$), 동차성($\|\alpha x\|=|\alpha|\|x\|$), 삼각부등식($\|x+y\|\le\|x\|+\|y\|$)을 만족하는 함수 $\|\cdot\|$이다. 이 책에서 가장 자주 쓰는 것은 $p\ge1$에 대한 $p$-노름
\[
\|x\|_p=\Bigl(\sum_{i=1}^n|x_i|^p\Bigr)^{1/p},\qquad\|x\|_\infty=\max_i|x_i|
\]
이며 특히 $\|x\|_2$를 Euclid 노름이라 한다.

\begin{proposition}[Hölder, Minkowski]
\label{prop:holder}
$1<p<\infty$, $1/p+1/q=1$이면 $|x^{\mathsf H}y|\le\|x\|_p\|y\|_q$이고, $p\ge1$이면 $\|x+y\|_p\le\|x\|_p+\|y\|_p$이다.
\end{proposition}

\begin{proof}
Young의 부등식 $ab\le a^p/p+b^q/q$($a,b\ge0$)는 $\ln$이 오목하다는 데서 나온다: $\ln(a^p/p+b^q/q)\ge(\ln a^p)/p+(\ln b^q)/q=\ln ab$. $x,y\ne0$이면 $a_i=|x_i|/\|x\|_p$, $b_i=|y_i|/\|y\|_q$에 Young을 쓰고 더하면 $\sum a_ib_i\le1/p+1/q=1$, 곧 Hölder이다. Minkowski는 $p>1$에서 $\sum|x_i+y_i|^p\le\sum|x_i||x_i+y_i|^{p-1}+\sum|y_i||x_i+y_i|^{p-1}$의 두 합에 Hölder를 쓰면 $(\|x\|_p+\|y\|_p)\|x+y\|_p^{p-1}$ 이하이므로 나온다($(p-1)q=p$). $p=1$은 항마다의 삼각부등식이다.
\end{proof}

유한차원에서는 어떤 노름을 쓰든 수렴이나 연속성 같은 질적인 성질이 같다. 다음 정리 때문이다.

\begin{theorem}[노름의 동치성]
\label{thm:norm-equiv}
$\mathbb{F}^n$ 위의 두 노름 $\|\cdot\|_\alpha$와 $\|\cdot\|_\beta$에 대해 $c\|x\|_\alpha\le\|x\|_\beta\le C\|x\|_\alpha$가 모든 $x$에서 성립하는 양의 상수 $c,C$가 있다. 특히
\[
\|x\|_\infty\le\|x\|_2\le\|x\|_1\le\sqrt n\,\|x\|_2\le n\|x\|_\infty
\]
이다.
\end{theorem}

\begin{proof}
$\|\cdot\|_\alpha$가 $\|\cdot\|_2$와 동치임을 보이면 충분하다. 표준기저 $e_i$로 $x=\sum x_ie_i$라 쓰면 삼각부등식과 Cauchy--Schwarz 부등식에서 $\|x\|_\alpha\le\sum|x_i|\|e_i\|_\alpha\le M\|x\|_2$, $M=(\sum\|e_i\|_\alpha^2)^{1/2}$이다. 이는 $\|\cdot\|_\alpha$가 $\|\cdot\|_2$에 대해 연속임을 뜻한다. 단위구면 $S=\{x:\|x\|_2=1\}$은 콤팩트하므로 $\|\cdot\|_\alpha$는 $S$ 위에서 최솟값 $m$을 갖고 $S$에 $0$이 없으므로 $m>0$이다. 동차성에서 모든 $x$에 대해 $m\|x\|_2\le\|x\|_\alpha\le M\|x\|_2$이다. 특정 부등식들은 $\max_i|x_i|^2\le\sum|x_i|^2\le(\sum|x_i|)^2$과 Cauchy--Schwarz 부등식 $\sum|x_i|\cdot1\le\sqrt n\|x\|_2$, 그리고 $\|x\|_2\le\sqrt n\|x\|_\infty$에서 나온다.
\end{proof}

\cref{thm:norm-equiv}의 상수에는 $n$이 들어 있다. 무한차원에서 이 정리가 깨지는 것도 그래서이다. 차원이 큰 계산에서는 노름을 무엇으로 고르느냐에 따라 오차 한계가 $\sqrt n$배나 $n$배까지 달라질 수 있다는 점도 기억해 두자.

\subsection{Computing the Euclidean Norm without Overflow}

$\|x\|_2$를 정의대로 계산하면 제곱하는 순간 문제가 생긴다. binary64가 담을 수 있는 가장 큰 수는 약 $1.8\times10^{308}=2^{1024}(1-2^{-53})$이므로 원소의 크기가 $2^{512}$을 넘으면 제곱이 넘쳐 $\infty$가 된다. 거꾸로 $2^{-511}$보다 작으면 제곱이 정규수의 아래 끝 $2^{-1022}$보다 작아져 비정규수가 되면서 정밀도를 잃고 $2^{-537.5}$ 이하이면 제곱이 $0$으로 반올림된다. 노름 자체는 원소와 같은 크기이므로 얼마든지 나타낼 수 있는데도 그렇다. 예컨대 numpy 2.5.5의 \code{np.linalg.norm}은 $(3\times10^{200},4\times10^{200})$에서 $\infty$, $(3\times10^{-200},4\times10^{-200})$에서 $0$을 준다(\code{figures/04/measure.py}).

이 문제는 척도를 바꾸어 피한다. $s=\|x\|_\infty$라 하면 $\|x\|_2=s\,\|x/s\|_2$이고 $x/s$의 원소는 모두 절댓값이 $1$ 이하이므로 제곱이 넘치지 않는다. 다만 나눗셈은 곱셈보다 비싸고 $s$를 구하려면 배열을 한 번 더 읽어야 한다. 대부분의 입력은 척도 없이 계산해도 값이 그대로 옳으므로 그런 입력을 가려 이 비용을 피해야 한다. 다음 명제가 그 판정이다. 아래에서 $u=2^{-53}$은 binary64의 반올림 단위이다(10장).

\begin{proposition}
\label{prop:norm-trust}
binary64 원소 $x_1,\dots,x_n$($n\ge1$)의 제곱을 반올림해 더한 합 $\hat s$가 유한하고 $\hat s\ge n\cdot2^{-1000}$이면 $\sqrt{\hat s}$는 합의 일반적인 오차 한계 안에서 $\|x\|_2$의 좋은 근사이다. 곧 넘침이 없고 비정규수 제곱의 절대오차들은 합쳐도 $\hat s$의 $2^{-75}$배를 넘지 않는다.
\end{proposition}

\begin{proof}
$\hat s$가 유한하면 어떤 제곱도, 어떤 중간 합도 넘치지 않았다(모두 음이 아니므로 중간 합은 최종 합 이하이다). 비정규수 영역의 반올림은 상대오차가 아니라 절대오차 $2^{-1075}$ 이하로 일어나므로 제곱 $n$개의 이런 오차는 $n\cdot2^{-1075}$ 이하이다. $\hat s\ge n\cdot2^{-1000}$이면 이 양은 $\hat s\cdot2^{-75}$ 이하이고 반올림 단위 $2^{-53}$보다 훨씬 작다. 나머지 반올림은 정규수 영역의 상대오차이므로 합의 오차 한계(11장)가 그대로 성립한다.
\end{proof}

stardust의 \code{vector\_norm}은 이 명제로 두 경로를 잇는다(\cref{alg:norm2-fast,alg:norm2}). 연속인 \code{float64}(복소수의 2-노름이면 실수부와 허수부를 두 항으로 본 \code{float64})는 먼저 제곱합만 빠르게 구해 명제로 판정한다. 판정을 통과하지 못한 입력과 그 밖의 dtype·배치는 최댓값과 NaN을 함께 모으는 조심스러운 경로로 간다.

\begin{algorithm}[htbp]
\begin{algorithmic}[1]
\Require 연속인 \code{float64} $x_1,\dots,x_n$ (메모리 순서)
\Ensure $\|x\|_2$, 또는 ``조심스러운 경로로''
\State $n$을 $2^{14}$항 덩어리로 자른다 \Comment{길이로만 정한 분할: 스레드 수와 무관}
\For{덩어리마다 (긴 합은 스레드로)}
  \State $512$항 이하: $16$갈래 \code{fma} 잎; $2048$항 이하: 그 잎의 쌍별 합; 더 길면 \code{cblas\_ddot}$(x,x)$
\EndFor
\State $\hat s\gets$ 덩어리 합들을 균형 트리로 더한 값
\State \textbf{if} $\hat s$가 유한하고 $\hat s\ge n\cdot2^{-1000}$ \textbf{then} \Return $\sqrt{\hat s}$ \Comment{\cref{prop:norm-trust}}
\State \Return 조심스러운 경로로 (\cref{alg:norm2}) \Comment{NaN, $\infty$, 넘침, 아주 작은 값}
\end{algorithmic}
\caption{2-노름의 빠른 경로 (stardust \code{fast\_norm}, \code{norm\_from\_squares})}
\label{alg:norm2-fast}
\end{algorithm}

\cref{alg:norm2-fast}는 길이에 따라 합하는 방식이 다르다. 짧은 합은 NEON 레지스터 여덟 개(갈래 $16$개)에 \code{fma}로 제곱을 쌓는 잎과 그 쌍별 합이고 $2048$항을 넘으면 덩어리마다 Accelerate의 \code{cblas\_ddot}(행렬 장치를 쓰고 누산기가 $256$개쯤이다)을 부른다. 덩어리 안에서는 누산기 여러 개로 나눈 합의 오차가, 덩어리 사이에서는 쌍별 합의 오차가 쌓인다. 누산기 $256$개에 $2^{14}$항을 나누면 한 누산기가 $64$항을 차례로 더하므로 오차는 쌍별 합의 $512$항 잎에서 갈래마다 $32$항을 더할 때와 같은 크기이다. 덩어리의 경계는 길이로만 정하므로 스레드 수가 결과 비트를 바꾸지 않는다. 덧셈 순서가 메모리 순서이므로 F 순서 행렬의 \code{'fro'}나 여러 축에 걸친 노름은 같은 값이라도 레이아웃에 따라 끝 비트가 다를 수 있다. 따로 빌드한 stardust에서 무작위 행렬 $20$개의 \code{'fro'}를 C 순서와 F 순서로 재면 $2$개가 달랐다. stardust는 속도를 얻으려고 설계에서 이 차이를 받아들였다.

\begin{algorithm}[htbp]
\begin{algorithmic}[1]
\Require 크기 $a_1,\dots,a_n$ (논리적 C 순서로 걷는다; 복소 2-노름은 실수부·허수부를 두 항으로)
\Ensure $\|a\|_2$
\State \textbf{if} $n=0$ \textbf{then} \Return $0$
\State $s\gets\sum_ia_i^2$ (쌍별 합, \code{fma}), $a_{\max}\gets\max_ia_i$, $\nu\gets$ NaN이 있는가 \Comment{한 번 읽기}
\State \textbf{if} $\nu$ \textbf{then} \Return NaN
\State \textbf{if} $a_{\max}=\infty$ \textbf{then} \Return $\infty$
\If{$s=\infty$ \textbf{or} ($a_{\max}>0$ \textbf{and} $\mathrm{fl}(a_{\max}^2)<2^{-969}$)} \Comment{넘침, 또는 가장 큰 제곱도 작다}
  \State $t\gets\sum_i(a_i/a_{\max})^2$ (쌍별 합) \Comment{넘침 없음, $t\ge1$}
  \State \Return $a_{\max}\sqrt t$
\EndIf
\State \Return $\sqrt s$
\end{algorithmic}
\caption{넘치지 않는 2-노름의 조심스러운 경로 (stardust \code{vnorm.hpp}의 \code{finish\_pnorm})}
\label{alg:norm2}
\end{algorithm}

\cref{alg:norm2}의 각 줄을 보자. 2행은 한 번의 읽기에서 세 가지를 모은다. 연속인 \code{float64}는 그 자리에서 읽으며 레지스터 여덟 개에 제곱을 쌓고 같은 원소의 절댓값으로 최댓값 레지스터 넷을 갱신한다(\code{sq\_leaf\_in\_place}). 연속이 아니거나 여러 축에 걸친 입력은 원소를 논리적 C 순서로 걸으며 $512$개씩 버퍼에 모아 같은 잎에 넘긴다(\code{Walk}). 축마다 남은 길이만큼 안쪽 축을 달리고 끝에 닿으면 바깥 축으로 올림을 넘기므로 힙 할당 없이 어떤 strides든 걷는다. 이 경로는 원소를 논리 순서로 더하므로 레이아웃과 무관한 비트를 준다. 3--4행은 특수값이다. NaN이 하나라도 있으면 NaN이고 NaN 없이 $\infty$가 있으면 $\infty$이다.

5행의 판정은 가장 큰 원소를 제곱한 값으로 한다. $a_{\max}$가 $2^{-537.5}$ 이하이면 그 제곱이 $0$으로 반올림되지만 $a_{\max}>0$을 따로 보므로 이 경우도 척도 맞춤으로 간다. 경계 $2^{-969}=2^{-1022}\cdot2^{53}$의 근거는 \cref{prop:norm-trust}와 같다. 가장 큰 제곱이 $2^{-969}$ 이상이면 합도 그 이상이고 비정규수 제곱의 절대오차 합 $n\cdot2^{-1075}$는 합의 $n\cdot2^{-106}$배 이하라 무시할 수 있다. 빠른 경로의 $\hat s<n\cdot2^{-1000}$은 같은 목적의 다른 판정이다. 6행은 $a_{\max}$로 나누므로 반올림이 한 번 더 든다. 넘침은 없고 아주 작은 $a_i/a_{\max}$가 비정규수가 되더라도 $a_{\max}$ 자신이 항 $1$을 주어 $t\ge1$이므로 그 오차는 무시할 만하다. 척도를 $a_{\max}$ 대신 그에 가까운 $2$의 거듭제곱으로 잡으면 나눗셈이 정확해진다(3장의 연분수 척도 맞춤과 같은 생각). stardust는 $a_{\max}$를 그대로 쓰는데 이 경로는 극단적인 입력에서만 타므로 영향이 작다. 그래서 드문 입력은 빠른 경로, 조심스러운 경로, 척도 맞춤에서 한 번씩 모두 세 번 읽힌다.

같은 문제를 한 번의 읽기로 푸는 방법도 있다. Blue의 방법\cite{blue1978}은 원소를 크기에 따라 큰 것, 보통인 것, 작은 것의 세 무리로 나누어 무리마다 다른 $2$의 거듭제곱 척도로 제곱을 누적하고 마지막에 합친다. 그러나 원소마다 분기가 들어 벡터화하기 어렵다. stardust는 흔한 경우를 벡터화한 한 번의 읽기로 끝내고 드문 경우에만 여러 번 읽는다.

\cref{fig:norm-scale}은 표준정규 난수 $1000$개를 $2^k$배한 벡터에서 세 방법의 상대오차를 잰 것이다. 참값은 제곱합을 정확한 유리수로 만들고 정수 제곱근으로 올바르게 반올림한 값이다. $2^k$를 곱하는 것은 정확하므로 넘침이나 아래넘침이 없는 한 세 방법이 모두 같은 값을 낸다. 제곱합이 약 $1000\approx2^{10}$이므로 순진한 계산과 numpy는 $k\ge510$에서 $\infty$를 내고 $k=-520$부터 비정규수로 정밀도를 잃어($k=-520$에서 $3.7\times10^{-14}$, $-530$에서 $3.1\times10^{-7}$) $k\le-540$에서는 $0$을 낸다. 따로 빌드한 stardust의 \code{vector\_norm}은 $k=-1070$부터 $1010$까지 모든 점에서 올바르게 반올림된 값을 냈다.

\begin{figure}[htbp]
\centering
\begin{tikzpicture}
\begin{semilogyaxis}[sdplot,xlabel={척도 지수 $k$},ylabel={상대오차},xmin=-1100,xmax=1040,ymin=1e-19,ymax=10,ytick={1e-18,1e-12,1e-6,1},yticklabels={${\le}10^{-18}$,$10^{-12}$,$10^{-6}$,$1$}]
\addplot[only marks,mark=*,mark size=1.3pt,black] table[x=k,y=naive]{figures/04/norm_scale.dat};
\addlegendentry{$\sqrt{\sum x_i^2}$}
\addplot[only marks,mark=o,mark size=2.2pt,black!50] table[x=k,y=numpy]{figures/04/norm_scale.dat};
\addlegendentry{numpy}
\addplot[only marks,mark=square,mark size=1.4pt,black] table[x=k,y=stardust]{figures/04/norm_scale.dat};
\addlegendentry{stardust}
\end{semilogyaxis}
\end{tikzpicture}
\caption{$2^k$배한 벡터의 2-노름의 상대오차(값 $1$은 $\infty$나 $0$ 같은 실패; 정확히 맞은 점은 $10^{-18}$에 그렸다).}
\label{fig:norm-scale}
\end{figure}

일반 $p$-노름과 다른 \code{ord}의 계산은 \cref{tab:vector-ords}에 모았다. $p$-노름은 $\sum a_i^p$를 \code{pow}로 쌓고 넘치거나 $a_{\max}^p<2^{-969}$이면 $a_{\max}(\sum(a_i/a_{\max})^p)^{1/p}$로 다시 계산한다. 이 척도 맞춤은 $p>0$에서만 한다. 음의 $p$는 큰 항이 작아지므로 넘침 걱정이 없다. 복소수에서 실수부와 허수부를 두 항으로 나누는 것은 2-노름뿐이고 다른 \code{ord}는 $|z|$를 \code{hypot}으로 구해 쓴다. 최댓값·최솟값·개수는 반올림이 없으므로 순서와 무관하다. 긴 벡터를 스레드로 나눌 때도 같은 덩어리나 같은 쌍별 트리의 잎을 나누어 맡기므로 결과 비트가 스레드 수와 무관하다(12장, 17장).

\begin{table}[htbp]
\begin{sdtable}{colspec={lX}, row{1}={font=\small\bfseries}}
\toprule
\code{ord} & 계산 \\
\midrule
$2$ & \cref{alg:norm2-fast}, 실패하면 \cref{alg:norm2}; 복소수는 두 항 \\
$1$ & $\sum|a_i|$; 연속 \code{float64}는 $1024$항까지 쌍별 합, 넘으면 \code{cblas\_dasum} \\
일반 $p$ & $(\sum a_i^p)^{1/p}$, $p>0$이면 넘침·아주 작은 값에서 척도 맞춤 \\
$\infty$, $-\infty$ & $\max|a_i|$, $\min|a_i|$ (NaN이면 NaN) \\
$0$ & $0$이 아닌 원소의 수 \\
\bottomrule
\end{sdtable}
\caption{\code{vector\_norm}의 \code{ord}별 계산(빈 입력은 모두 $0$).}
\label{tab:vector-ords}
\end{table}

\section{Matrix Norms}
\label{sec:matrix-norms}

\subsection{Induced Norms}

행렬을 길이 $mn$인 벡터로 보면 벡터 노름을 그대로 쓸 수 있지만 행렬은 벡터를 보내는 사상이므로 ``벡터를 얼마나 늘리는가''로 크기를 재는 편이 더 자연스럽다.

\begin{definition}[유도 노름]
\label{def:induced}
$\mathbb{F}^n$과 $\mathbb{F}^m$의 벡터 노름에 대해 $m\times n$ 행렬 $A$의 \term{유도 노름}{induced norm}을
\[
\|A\|=\max_{x\ne0}\frac{\|Ax\|}{\|x\|}=\max_{\|x\|=1}\|Ax\|
\]
로 정의한다. 벡터 노름이 $p$-노름이면 $\|A\|_p$로 쓴다.
\end{definition}

정의에서 $\|Ax\|\le\|A\|\|x\|$이고 이를 두 번 쓰면 \term{부분곱셈성}{submultiplicativity} $\|AB\|\le\|A\|\|B\|$를 얻는다. 단위구면이 콤팩트하고 $x\mapsto\|Ax\|$가 연속이므로 최댓값이 존재한다. 유도 노름 가운데 셋은 원소에서 바로 계산된다.

\begin{theorem}
\label{thm:induced-formulas}
$m\times n$ 행렬 $A$에 대해 다음이 성립한다.
\begin{enumerate}
\item $\|A\|_1=\max_j\sum_i|a_{ij}|$, 곧 열의 절댓값 합의 최댓값이다.
\item $\|A\|_\infty=\max_i\sum_j|a_{ij}|$, 곧 행의 절댓값 합의 최댓값이다.
\item $\|A\|_2$는 $A^{\mathsf H}A$의 가장 큰 고유값의 제곱근, 곧 $A$의 가장 큰 특이값이다.
\end{enumerate}
\end{theorem}

\begin{proof}
(i) $\|x\|_1=1$이면 $\|Ax\|_1=\sum_i|\sum_ja_{ij}x_j|\le\sum_j|x_j|\sum_i|a_{ij}|\le\max_j\sum_i|a_{ij}|$이고 등호는 $x$를 그 최댓값을 주는 열의 표준기저벡터로 잡으면 성립한다. (ii) $\|x\|_\infty=1$이면 각 행에서 $|\sum_ja_{ij}x_j|\le\sum_j|a_{ij}|$이고 등호는 최댓값을 주는 행 $i$에 대해 $x_j=\overline{a_{ij}}/|a_{ij}|$($a_{ij}=0$이면 $x_j=1$)로 잡으면 성립한다. (iii) $\|Ax\|_2^2=x^{\mathsf H}(A^{\mathsf H}A)x$이고 $A^{\mathsf H}A$는 에르미트이며 양의 준정부호이므로 고유값 $\lambda_1\ge\cdots\ge\lambda_n\ge0$과 정규직교 고유벡터 $v_i$가 있다. $x=\sum c_iv_i$, $\sum|c_i|^2=1$이면 $\|Ax\|_2^2=\sum\lambda_i|c_i|^2\le\lambda_1$이고 $x=v_1$에서 등호가 성립한다.
\end{proof}

\cref{thm:induced-formulas}의 (i)과 (ii)는 원소를 한 번 훑으면 계산되지만 (iii)은 특이값을 구해야 한다. 그래서 큰 행렬에서는 $\|A\|_2$ 대신 $\|A\|_1$이나 $\|A\|_\infty$를 쓰거나 $\|A\|_2$를 어림하는 경우가 많다. 이들 사이에는 \cref{thm:norm-equiv}에서 나오는 차원 상수의 관계 외에 $\|A\|_2\le\sqrt{\|A\|_1\|A\|_\infty}$도 성립한다. $\|A\|_2^2=\rho(A^{\mathsf H}A)\le\|A^{\mathsf H}A\|_1\le\|A^{\mathsf H}\|_1\|A\|_1=\|A\|_\infty\|A\|_1$이기 때문이다. 여기서 첫 부등식은 스펙트럼 반지름이 모든 유도 노름 이하라는 사실이다. 고유값 $\lambda$와 고유벡터 $v$에 대해 $|\lambda|\|v\|=\|Mv\|\le\|M\|\|v\|$이기 때문이다.

\subsection{The Frobenius Norm and Unitary Invariance}

행렬을 벡터로 본 Euclid 노름 $\|A\|_F=(\sum_{i,j}|a_{ij}|^2)^{1/2}=(\operatorname{tr}A^{\mathsf H}A)^{1/2}$을 \term{Frobenius 노름}{Frobenius norm}이라 한다. 이 노름은 유도 노름이 아니지만(단위행렬의 Frobenius 노름이 $\sqrt n$이므로) 부분곱셈적이며 $\|A\|_2\le\|A\|_F\le\sqrt{\operatorname{rank}A}\,\|A\|_2$를 만족한다. 부분곱셈성은 $AB$의 $(i,j)$ 원소에 Cauchy--Schwarz를 쓴 $|\sum_ka_{ik}b_{kj}|^2\le(\sum_k|a_{ik}|^2)(\sum_k|b_{kj}|^2)$를 모든 $i,j$에 대해 더하면 나온다. 두 부등식은 다음 절의 특이값 분해로 보면 $\|A\|_F^2=\sum_{i\le r}\sigma_i^2$이 $\sigma_1^2$ 이상, $r\sigma_1^2$ 이하이므로 성립한다. $\|\cdot\|_2$와 $\|\cdot\|_F$는 모두 유니터리 행렬 $U$, $V$에 대해 $\|UAV\|=\|A\|$인 \term{유니터리 불변}{unitarily invariant} 노름이다. 유니터리 변환은 Euclid 길이를 바꾸지 않으므로 $\|UAV\|_2=\|A\|_2$이고 $\operatorname{tr}((UAV)^{\mathsf H}UAV)=\operatorname{tr}(V^{\mathsf H}A^{\mathsf H}AV)=\operatorname{tr}(A^{\mathsf H}A)$이므로 $\|UAV\|_F=\|A\|_F$이다. 다음 절에서 이 성질을 특이값 분해와 함께 쓰면 두 노름을 특이값만으로 나타낼 수 있다.

\subsection{Computing Matrix Norms}

stardust의 \code{matrix\_norm}은 \code{ord}에 따라 세 갈래로 나뉜다(\code{csrc/linalg/ops/norms.cpp}). \code{'fro'}는 행렬을 벡터로 보고 \cref{alg:norm2}를 쓴다. $2$, $-2$, \code{'nuc'}은 특이값을 구해 가장 큰 값, 가장 작은 값, 합을 돌려준다. 특이값은 LAPACK을 부르므로 NaN이나 $\infty$가 있으면 오류를 내고 다른 \code{ord}는 NaN을 그대로 퍼뜨린다. 남은 $\pm1$과 $\pm\infty$가 \cref{thm:induced-formulas}의 선 합이다(\cref{alg:line-norm}).

\begin{algorithm}[htbp]
\begin{algorithmic}[1]
\Require $m\times n$ 행렬 $A$ (shape, 바이트 strides), \code{ord}$\in\{1,-1,\infty,-\infty\}$
\Ensure 선 합 $\sum|a|$의 최댓값(\code{ord}$>0$) 또는 최솟값
\State 선 $\gets$ (\code{ord}$=\pm1$이면 열, 아니면 행); $L\gets$ 선의 수 \Comment{$L=0$이면 $0$}
\If{선 안의 원소가 메모리에서 연속}
  \State 선마다 $\sum|a|$를 그 자리에서 쌍별로 \Comment{한 선이 한 번의 빠른 합}
\ElsIf{선들이 메모리에서 나란하다 (선 사이 간격이 원소 하나)}
  \State 선 $w$개씩 띠로 묶어 행을 차례로 읽으며 $w$개의 합을 레지스터에 쌓는다 \Comment{띠 안은 반으로 나누어 더함}
\Else
  \State 선마다 원소를 논리 순서로 모아 쌍별 합
\EndIf
\State \Return NaN이 있으면 NaN, 아니면 합들의 최댓값(최솟값)
\end{algorithmic}
\caption{$\pm1$, $\pm\infty$ 노름의 선 합 (stardust \code{matrix\_norm})}
\label{alg:line-norm}
\end{algorithm}

\cref{alg:line-norm}의 세 갈래는 메모리 배치에서 나온다. C 순서 행렬의 $\infty$-노름은 행 합이므로 행 하나가 연속이고 2행의 갈래를 탄다. 같은 행렬의 $1$-노름은 열 합인데 열 하나의 원소들은 행 길이만큼 떨어져 있다. 열마다 따로 더하면 원소 하나를 읽을 때마다 다른 캐시 줄을 건드리므로 느리다. 대신 열들이 메모리에서 나란하다는 점을 써서 열 $w$개를 띠로 묶고 행을 하나씩 내려가며 $w$개의 합을 동시에 쌓는다(4행). 행 하나 안의 $w$개 원소는 연속이므로 벡터 명령으로 읽힌다. stardust는 열 $16$개를 레지스터 여덟 개에 두고 띠의 폭을 $128$로, 병렬일 때는 스레드마다 일이 돌아가도록 $32$로 잡는다. 띠 안의 행 방향 합도 $32$행 잎 위로는 반씩 나누어 더한다. 그래서 한 합에서 어떤 항이 거치는 덧셈은 잎 안의 차례 덧셈과 트리의 깊이를 합쳐 $64+\log_2n$ 남짓이고 stardust의 설계 문서는 이를 바탕으로 띠 합의 오차를 $2(64+\log_2n+2)\,u\sum|a|$로 누른다(합의 일반적인 한계는 11장). stardust의 시험은 이 한계로 띠 합을 \code{sum}의 결과와 견준다.

다만 같은 값의 행렬이라도 C 순서이면 띠로, F 순서이면 선마다 따로 더하므로 덧셈의 순서가 달라 끝 비트가 다를 수 있다. 무작위 행렬 $40$개($n=100$--$3000$)의 $1$-노름과 $\infty$-노름을 C 순서와 F 순서로 재어 보면 $17$개가 마지막 자리에서 달랐다. stardust는 설계 단계에서 띠 방식이 C 순서 $1000\times1000$의 $1$-노름을 약 세 배 빠르게 한다는 점을 보고 이 레이아웃 의존을 받아들였다. 대신 \code{cond}처럼 노름을 내부에서 쓰는 함수는 패킹한 열 우선 버퍼에서 계산해 레이아웃과 무관하다. \cref{tab:matrix-norm-time}은 $4096\times4096$ 행렬에서 잰 시간이다.

\begin{table}[htbp]
\begin{sdtable}{colspec={lccc}, row{1}={font=\small\bfseries}}
\toprule
\code{ord} & stardust, C 순서 & stardust, 전치 뷰 & numpy, C 순서 \\
\midrule
$1$ & $1.04$ & $0.71$ & $4.03$ \\
$\infty$ & $0.72$ & $1.05$ & $3.94$ \\
\code{'fro'} & $0.53$ & $0.53$ & $0.98$ \\
\bottomrule
\end{sdtable}
\caption{$4096\times4096$ \code{float64} 행렬 노름의 시간(ms, Apple M4 Pro, 기본 스레드, 다섯 번 중 최솟값, \code{figures/04/measure.py}). 전치 뷰는 같은 메모리를 F 순서로 본 것이다.}
\label{tab:matrix-norm-time}
\end{table}

표에서 C 순서의 $\infty$-노름(행이 연속)과 전치 뷰의 $1$-노름(열이 연속)은 같은 메모리를 같은 방식으로 읽으므로 시간이 비슷하다. 띠를 타는 나머지 두 경우는 조금 느리다. 어느 쪽이든 원소를 한 번씩만 읽는다.

\section{The Singular Value Decomposition}
\label{sec:svd}

\subsection{Existence}

\begin{theorem}[특이값 분해]
\label{thm:svd}
$m\times n$ 행렬 $A$에 대해 유니터리 행렬 $U\in\mathbb{F}^{m\times m}$, $V\in\mathbb{F}^{n\times n}$와 대각 성분이 $\sigma_1\ge\sigma_2\ge\cdots\ge\sigma_{\min(m,n)}\ge0$인 $m\times n$ 대각행렬 $\Sigma$가 있어 $A=U\Sigma V^{\mathsf H}$이다. $\sigma_i$를 $A$의 \term{특이값}{singular value}이라 하며 특이값은 $A$로 유일하게 정해진다.
\end{theorem}

\begin{proof}
$\min(m,n)$에 대한 귀납법을 쓴다. $A=0$이면 자명하므로 $A\ne0$이라 하자. $\sigma_1=\|A\|_2$로 두고 $\|v_1\|_2=1$, $\|Av_1\|_2=\sigma_1$인 $v_1$을 잡은 뒤 $u_1=Av_1/\sigma_1$로 두자. $v_1$과 $u_1$을 각각 유니터리 행렬 $V_1=[v_1\ \widetilde V]$, $U_1=[u_1\ \widetilde U]$로 넓히면
\[
U_1^{\mathsf H}AV_1=\begin{bmatrix}\sigma_1&w^{\mathsf H}\\0&B\end{bmatrix}
\]
꼴이다. 첫 열이 $(\sigma_1,0,\dots,0)^{\mathsf T}$인 것은 $Av_1=\sigma_1u_1$이고 $\widetilde U^{\mathsf H}u_1=0$이기 때문이다. 이제 $w=0$임을 보이자. 이 행렬을 $S$라 하면 $\|S\|_2=\|A\|_2=\sigma_1$이고
\[
\Bigl\|S\begin{bmatrix}\sigma_1\\w\end{bmatrix}\Bigr\|_2\ge\sigma_1^2+w^{\mathsf H}w=(\sigma_1^2+\|w\|_2^2)^{1/2}\Bigl\|\begin{bmatrix}\sigma_1\\w\end{bmatrix}\Bigr\|_2
\]
이다. 왼쪽 부등호는 $S(\sigma_1,w)^{\mathsf T}$의 첫 성분이 $\sigma_1^2+w^{\mathsf H}w$라는 데서 나온다. 따라서 $\|S\|_2\ge(\sigma_1^2+\|w\|_2^2)^{1/2}$인데 $\|S\|_2=\sigma_1$이므로 $w=0$이다. 남은 블록 $B$에 귀납 가정을 쓰면 분해를 얻고 $\|B\|_2\le\|S\|_2=\sigma_1$이므로 특이값의 순서도 맞다. 유일성은 $A^{\mathsf H}A=V\Sigma^{\mathsf H}\Sigma V^{\mathsf H}$에서 $\sigma_i^2$이 $A^{\mathsf H}A$의 고유값이라는 데서 나온다.
\end{proof}

특이값 분해를 알면 행렬에 관한 정보를 거의 다 읽어 낼 수 있다. $\sigma_r>0=\sigma_{r+1}$이면 $A$의 랭크는 $r$이고 $U$의 처음 $r$개 열은 치역의, $V$의 나머지 열은 영공간의 정규직교기저이다. 유니터리 불변성에서 $\|A\|_2=\sigma_1$이고 $\|A\|_F=(\sum\sigma_i^2)^{1/2}$이며 특이값의 합 $\|A\|_*=\sum\sigma_i$를 \term{핵 노름}{nuclear norm}이라 한다. 핵 노름의 삼각부등식은 쌍대 표현 $\|A\|_*=\max\{\operatorname{Re}\operatorname{tr}(B^{\mathsf H}A):\|B\|_2\le1\}$에서 나온다. 최댓값은 $B=UV^{\mathsf H}$(얇은 분해)에서 이루어지고 임의의 $\|B\|_2\le1$에 대해 $\operatorname{Re}\operatorname{tr}(B^{\mathsf H}A)=\sum_i\sigma_i\operatorname{Re}(u_i^{\mathsf H}Bv_i)\le\sum\sigma_i$이다. 이처럼 각 $B$에 대한 일차식들의 최댓값으로 나타낸 함수는 삼각부등식을 만족한다. 정사각 가역행렬이면 $\|A^{-1}\|_2=1/\sigma_n$이다. 특이값을 계산하는 알고리즘은 56장에서 다룬다.

\subsection{Best Low-Rank Approximation}

특이값 분해는 여러 응용 가운데에서도 행렬을 랭크가 낮은 행렬로 근사하는 문제에 정확한 답을 주는 것으로 가장 유명하다.

\begin{theorem}[Eckart--Young]
\label{thm:eckart-young}
$A=\sum_{i}\sigma_iu_iv_i^{\mathsf H}$를 특이값 분해라 하고 $k<\operatorname{rank}A$에 대해 $A_k=\sum_{i=1}^k\sigma_iu_iv_i^{\mathsf H}$로 두면
\[
\min_{\operatorname{rank}B\le k}\|A-B\|_2=\|A-A_k\|_2=\sigma_{k+1}
\]
이다.
\end{theorem}

\begin{proof}
$A-A_k=\sum_{i>k}\sigma_iu_iv_i^{\mathsf H}$의 가장 큰 특이값은 $\sigma_{k+1}$이므로 등식의 오른쪽 부분은 분명하다. 이제 랭크가 $k$ 이하인 임의의 $B$에 대해 $\|A-B\|_2\ge\sigma_{k+1}$임을 보이자. $B$의 영공간은 차원이 $n-k$ 이상이고 $v_1,\dots,v_{k+1}$이 생성하는 공간은 차원이 $k+1$이므로 두 공간의 차원의 합이 $n$을 넘어 $0$이 아닌 공통 원소 $z$가 있다. $\|z\|_2=1$로 잡으면 $Bz=0$이므로
\[
\|A-B\|_2^2\ge\|(A-B)z\|_2^2=\|Az\|_2^2=\sum_{i=1}^{k+1}\sigma_i^2|v_i^{\mathsf H}z|^2\ge\sigma_{k+1}^2
\]
이다. 마지막 부등호는 $\sum_{i\le k+1}|v_i^{\mathsf H}z|^2=\|z\|_2^2=1$에서 나온다.
\end{proof}

이 정리는 Eckart와 Young이 1936년 심리학 학술지에 Frobenius 노름의 경우로 발표했다\cite{eckart1936young}. 같은 $A_k$가 Frobenius 노름에서도 최적이며 그때의 오차는 $(\sum_{i>k}\sigma_i^2)^{1/2}$이다. 증명은 Weyl 부등식의 한 꼴 $\sigma_{i+k}(A)\le\sigma_i(A-B)+\sigma_{k+1}(B)$(\cref{thm:courant-fischer}에서 부분공간의 차원을 세어 얻는다)에서 출발한다. 여기에 $\operatorname{rank}B\le k$, 곧 $\sigma_{k+1}(B)=0$을 넣으면 $\sigma_i(A-B)\ge\sigma_{i+k}(A)$이고 제곱해 더하면 $\|A-B\|_F^2\ge\sum_{i>k}\sigma_i^2$이다. 주성분 분석, 잡음 제거, 추천 시스템의 행렬 완성이 모두 이 정리에 기댄다.

\subsection{Numerical Rank and the Pseudoinverse}

\cref{thm:eckart-young}을 뒤집어 읽으면 $\sigma_{k+1}$은 $A$와 ``랭크가 $k$ 이하인 가장 가까운 행렬'' 사이의 거리이다. 그런데 반올림된 행렬에서는 정확한 $0$인 특이값이 거의 나오지 않으므로 랭크를 ``$0$이 아닌 특이값의 개수''로 정의하면 쓸모가 없다. 대신 허용오차 $\tau$를 정해 $\sigma_i>\tau$인 특이값의 개수, 곧 ``$\tau$ 안에서 $A$와 구별할 수 없는 행렬들 가운데 가장 낮은 랭크''를 \term{수치적 랭크}{numerical rank}로 삼는다. 같은 생각으로 \term{의사역행렬}{pseudoinverse}을 정의한다. $A=U\Sigma V^{\mathsf H}$에서 $\Sigma^+$를 $\Sigma$의 전치에서 $\tau$보다 큰 대각 성분을 그 역수로, 나머지를 $0$으로 바꾼 행렬이라 하면 $A^+=V\Sigma^+U^{\mathsf H}$이다. $\tau$가 $0$이 아닌 가장 작은 특이값보다 작을 때 $x=A^+b$가 최소제곱 문제 $\min\|Ax-b\|_2$의 해 가운데 노름이 가장 작은 해라는 사실은 52장에서 증명한다.

\begin{algorithm}[htbp]
\begin{algorithmic}[1]
\Require $m\times n$ 행렬 $A$, 상대 허용오차 \code{rtol} (없으면 $\max(m,n,1)\cdot2^{-52}$)
\Ensure 수치적 랭크 $r$, 또는 의사역행렬 $A^+$
\State $\sigma_1\ge\dots\ge\sigma_k\gets$ $A$의 특이값 ($k=\min(m,n)$) \Comment{랭크만이면 $U$, $V$ 없이}
\State $\tau\gets\code{rtol}\cdot\sigma_1$ ($k=0$이면 $0$)
\State $r\gets\#\{i:\sigma_i>\tau\}$ \Comment{$\code{rtol}\ge1$이면 $r=0$}
\State (\code{pinv}) $W\gets V\operatorname{diag}(\rho_i)$, $\rho_i=1/\sigma_i$ ($\sigma_i>\tau$이고 $\sigma_i\ge2^{-1022}$), 아니면 $0$ \Comment{$n\times k$}
\State (\code{pinv}) $A^+\gets WU^{\mathsf H}$ (\code{gemm} 한 번)
\State (\code{pinv}) 남긴 $\sigma_i$ 가운데 $2^{-1022}$보다 작은 것이 있으면 $2^{600}$배 이동해 따로 곱해 더한다
\end{algorithmic}
\caption{수치적 랭크와 의사역행렬 (stardust \code{matrix\_rank}, \code{pinv})}
\label{alg:rank-pinv}
\end{algorithm}

\cref{alg:rank-pinv}의 1행은 \code{matrix\_rank}에서는 특이값만 구하고 \code{pinv}에서는 얇은 분해 $A=U_k\Sigma_kV_k^{\mathsf H}$(LAPACK \code{gesdd}, $U$는 $m\times k$)까지 구한다. NaN이나 $\infty$가 있으면 LAPACK에 넘기기 전에 오류를 내고 빈 행렬이면 랭크 $0$과 $n\times m$ 영행렬을 돌려준다. 4--5행은 버린 특이값의 역수를 $0$으로 두고 $k$개 열 전부를 \code{gemm} 한 번에 넘긴다. 남긴 열만 골라 모으는 것보다 단순하고 비용 차이가 작다. 6행은 드문 경우를 처리한다. 남긴 특이값이 정규수의 아래 끝 $2^{-1022}$보다 작으면 역수가 $2^{1022}$을 넘어 넘칠 수 있으므로 그 특이값에 $2^{600}$을 곱한 뒤 역수를 만들어 따로 곱하고 결과에 $2^{600}$을 곱해 더한다. $2$의 거듭제곱 척도이므로 반올림이 들지 않는다. 비용은 특이값 분해의 $O(mn\min(m,n))$에 \code{gemm}의 $2mnk$이다.

허용오차의 기본값 $\max(m,n)\cdot2^{-52}\cdot\sigma_1$은 numpy의 \code{matrix\_rank}와 같고 Array API 표준의 \code{pinv}와 같다(\code{csrc/linalg/args.cpp}의 \code{parse\_rtol}). numpy의 \code{pinv}는 따로 정하지 않으면 $10^{-15}\cdot\sigma_1$을 쓰므로 그 경우와는 다르다. 이 값은 \cref{sec:perturbation}의 Weyl 부등식에서 나온다. 역방향으로 안정한 특이값 알고리즘은 $\|E\|_2$가 $u\|A\|_2$의 작은 배수인 $A+E$의 특이값을 정확히 계산하므로 특이값의 절대오차가 $\sigma_1$에 비례하는 크기이다. 그 크기 아래의 특이값은 참값이 $0$인지 아닌지 계산으로 가릴 수 없다. 차원 $\max(m,n)$은 그 배수가 차원과 함께 자란다는 어림이다. \code{rtol}은 음이 아닌 유한한 수여야 하고 $1$ 이상을 주면 모든 특이값이 $0$으로 세어진다.

\cref{fig:svd-err}로 이 사정을 확인해 보자. 정수 벡터로 만든 Householder 반사 $I-2vv^{\mathsf T}/(v^{\mathsf T}v)$는 유리수 직교행렬이므로 $Q_1\operatorname{diag}(\sigma_i)Q_2$를 정확한 유리수로 만들 수 있다. 특이값을 $\sigma_i=2^{-k_i}$($k_i$는 $52(i-1)/39$를 반올림한 서로 다른 정수 $40$개)로 정하고 이 행렬을 binary64로 반올림해 stardust의 \code{svdvals}로 계산하면 상대오차가 $u\sigma_1/\sigma_i$의 선보다 두세 자리 아래에서 거의 같은 기울기로 자란다. 가장 작은 $\sigma_{40}=2^{-52}$는 $2.2208\times10^{-16}$으로 나와 상대오차가 $1.7\times10^{-4}$이다. 행렬을 반올림하는 순간 이미 $u\sigma_1$ 크기의 섭동이 들어가므로 이 오차는 알고리즘의 결함으로 볼 수 없고 문제 자체에서 온다.

\begin{figure}[htbp]
\centering
\begin{tikzpicture}
\begin{semilogyaxis}[sdplot,xlabel={특이값의 번호 $i$},ylabel={$\sigma_i$의 상대오차},xmin=0,xmax=41,ymin=1e-18,ymax=1e-2]
\addplot[only marks,mark=*,mark size=1.3pt,black] table[x=i,y=relerr]{figures/04/svd_err.dat};
\addlegendentry{stardust \code{svdvals}}
\addplot[black!50,dashed] table[x=i,y=bound]{figures/04/svd_err.dat};
\addlegendentry{$u\sigma_1/\sigma_i$}
\end{semilogyaxis}
\end{tikzpicture}
\caption{특이값을 $1$부터 $2^{-52}$까지 정해 만든 $40\times40$ 행렬의 \code{svdvals} 상대오차(\code{figures/04/measure.py}; $10^{-18}$은 정확히 맞은 점).}
\label{fig:svd-err}
\end{figure}


\section{Condition Numbers}
\label{sec:condition}

\subsection{The Condition Number of a Problem}

조건수는 ``입력이 조금 바뀔 때 답이 얼마나 바뀌는가''를 재는 수이며 알고리즘이 아니라 \textbf{문제}의 성질이다. 미분가능한 $f$로 표현되는 문제에서 입력 $x$의 상대 섭동이 답에 주는 상대 변화의 비율을 \term{상대 조건수}{relative condition number}라 한다. 스칼라 함수라면 $\kappa_f(x)=|xf'(x)/f(x)|$이고 11장에서는 이 값으로 오차가 어디서 커지는지 판단한다. 여기서는 선형 연립방정식의 조건수를 자세히 본다.

\begin{definition}[행렬의 조건수]
\label{def:cond}
가역행렬 $A$와 행렬 노름 $\|\cdot\|$에 대해 $\kappa(A)=\|A\|\|A^{-1}\|$을 $A$의 \term{조건수}{condition number}라 한다. 2-노름에서는 $\kappa_2(A)=\sigma_1/\sigma_n$이다.
\end{definition}

\begin{theorem}
\label{thm:linsys-perturb}
가역행렬 $A$와 $b\ne0$에 대해 $Ax=b$라 하자. $\|A^{-1}\|\|\Delta A\|<1$이고 $(A+\Delta A)(x+\Delta x)=b+\Delta b$이면
\[
\frac{\|\Delta x\|}{\|x\|}\le\frac{\kappa(A)}{1-\kappa(A)\|\Delta A\|/\|A\|}\Bigl(\frac{\|\Delta A\|}{\|A\|}+\frac{\|\Delta b\|}{\|b\|}\Bigr)
\]
이다.
\end{theorem}

\begin{proof}
두 식을 빼면 $A\Delta x=\Delta b-\Delta A(x+\Delta x)$이므로 $\Delta x=A^{-1}(\Delta b-\Delta Ax-\Delta A\Delta x)$이다. 노름을 취하면
\[
\|\Delta x\|\le\|A^{-1}\|\bigl(\|\Delta b\|+\|\Delta A\|\|x\|\bigr)+\|A^{-1}\|\|\Delta A\|\|\Delta x\|
\]
이고 $\|A^{-1}\|\|\Delta A\|<1$이므로 $\|\Delta x\|$를 정리할 수 있다. 양변을 $\|x\|$로 나누고 $\|b\|=\|Ax\|\le\|A\|\|x\|$, 곧 $1/\|x\|\le\|A\|/\|b\|$를 쓰면 주장이 나온다.
\end{proof}

\cref{thm:linsys-perturb}에 따르면 입력의 상대오차가 $\epsilon$이면 해의 상대오차는 대략 $\kappa(A)\epsilon$까지 커질 수 있다. 반올림된 입력의 상대오차는 이미 $u\approx10^{-16}$ 정도이므로 $\kappa(A)\approx10^{k}$이면 어떤 알고리즘을 쓰더라도 해의 유효숫자를 대략 $k$자리 잃을 각오를 해야 한다. 한편 48장의 부분 피벗 LU처럼 역방향으로 안정한 알고리즘(LU는 성장 인자가 작은 한)은 $(A+\Delta A)\hat x=b$, $\|\Delta A\|\lesssim c_nu\|A\|$를 만족하는 $\hat x$를 주므로 정리에 넣으면 전방 오차가 $\kappa(A)c_nu$ 정도로 눌린다. 이런 알고리즘의 오차는 문제가 허락하는 수준을 넘지 않는다.

이 예측을 Hilbert 행렬 $h_{ij}=1/(i+j-1)$로 확인하자. $H$를 binary64로 만들고 $b=H\mathbf1$을 binary64로 계산한 뒤 stardust의 \code{solve}로 풀고 같은 binary64 행렬과 우변의 연립방정식을 정확한 유리수 Gauss 소거로 풀어 참값으로 삼았다. 이렇게 하면 행렬을 만드는 반올림과 무관하게 \code{solve} 자신의 오차만 남는다. \cref{fig:hilbert}에서 상대오차는 $\kappa_2u$의 선을 따라 늘어나 $n=8$($\kappa_2=1.5\times10^{10}$)에서 $5\times10^{-8}$, $n=12$($\kappa_2=1.7\times10^{16}$)에서 $10^{-2}$이다. $\kappa_2u$가 $1$을 넘는 $n=12$부터 오차가 커지고 $\kappa_2$가 $1/u$의 수백 배인 $n=13$에서는 답이 거의 남지 않는다. 그림의 $\kappa_2$는 \code{cond}로 구하지 않았다. 계산한 $\sigma_n$은 \cref{fig:svd-err}처럼 절대오차가 $u\sigma_1$ 크기이므로 $\kappa_2$가 $1/u$에 가까우면 믿을 수 없다. 대신 가장 큰 특이값은 상대적으로 정확하다는 점을 써서 binary64 $H$의 정확한 유리수 역행렬을 binary64로 반올림하고 두 행렬의 가장 큰 특이값을 곱했다(\code{figures/04/hilbert\_kappa.py}).

\begin{figure}[htbp]
\centering
\begin{tikzpicture}
\begin{loglogaxis}[sdplot,xlabel={$\kappa_2(H)$},ylabel={해의 상대오차},xmin=5,xmax=1e19,ymin=1e-17,ymax=10]
\addplot[only marks,mark=*,mark size=1.4pt,black] table[x=kappa,y=err]{figures/04/hilbert.dat};
\addlegendentry{\code{solve}}
\addplot[black!50,dashed,domain=5:1e19,samples=2] {1.11e-16*x};
\addlegendentry{$\kappa_2u$}
\end{loglogaxis}
\end{tikzpicture}
\caption{$n=2,\dots,13$의 Hilbert 행렬로 푼 연립방정식의 상대오차(stardust \code{solve}, 참값은 같은 binary64 입력의 정확한 유리수 해).}
\label{fig:hilbert}
\end{figure}

조건수는 거리로도 읽힌다. 그 역수가 가역행렬과 가장 가까운 특이행렬 사이의 상대 거리이다.

\begin{theorem}
\label{thm:dist-singular}
가역행렬 $A$에 대해 $\min\{\|\Delta A\|_2:A+\Delta A\text{가 특이}\}=\sigma_n=\|A\|_2/\kappa_2(A)$이다.
\end{theorem}

\begin{proof}
\cref{thm:eckart-young}을 $k=n-1$에 적용하면 랭크가 $n-1$ 이하인 행렬 $B$까지의 거리 $\|A-B\|_2$의 최솟값은 $\sigma_n$이다. 특이행렬은 랭크가 $n-1$ 이하인 행렬과 같으므로 주장이 나온다.
\end{proof}

$n=13$의 Hilbert 행렬은 \cref{thm:dist-singular}로 보면 상대 거리 $1/\kappa_2\approx4.6\times10^{-19}$ 안에 특이행렬이 있다. 반올림 단위보다 가까우므로 binary64로 저장한 순간 그 행렬은 특이행렬과 구별되지 않는다. stardust의 \code{cond}는 \code{ord=2}이면 \cref{def:cond} 그대로 특이값으로 $\sigma_1/\sigma_n$을 계산하고 $\sigma_n=0$이면 $\infty$를 돌려준다. \code{ord=-2}는 $\sigma_n/\sigma_1$이고 빈 행렬은 오류이다. 다른 \code{ord}는 $\|A\|\|A^{-1}\|$을 LU로 만든 역행렬에서 구하는데 LU가 특이를 알리거나 역행렬이 넘치면 $\infty$를 돌려준다. 이 노름은 패킹한 열 우선 버퍼에서 계산하므로 레이아웃과 무관하다. 그러나 앞의 문단이 보여 주듯 계산한 $\sigma_n$의 절대오차가 $u\sigma_1$이므로 $\kappa_2$가 $1/u$에 가까워지면 \code{cond} 자신도 크기만 맞는다. 큰 행렬의 조건수를 싸게 어림하는 알고리즘은 53장에서 다룬다.

\section{Perturbation of Eigenvalues and Singular Values}
\label{sec:perturbation}

\subsection{The Minimax Principle and Weyl's Inequality}

에르미트 행렬의 고유값과 모든 행렬의 특이값은 섭동에 아주 안정적이다. 이 안정성은 고유값을 최적화 문제의 값으로 나타내는 다음 정리에서 나온다.

\begin{theorem}[Courant--Fischer]
\label{thm:courant-fischer}
$n\times n$ 에르미트 행렬 $H$의 고유값을 $\lambda_1\ge\cdots\ge\lambda_n$이라 하면
\[
\lambda_k=\max_{\dim S=k}\ \min_{x\in S,\ \|x\|_2=1}x^{\mathsf H}Hx
\]
이다. 여기서 최댓값은 $\mathbb{F}^n$의 $k$차원 부분공간 $S$에 대해 잡는다.
\end{theorem}

\begin{proof}
정규직교 고유벡터 $v_1,\dots,v_n$을 잡자. $S=\operatorname{span}(v_1,\dots,v_k)$로 두면 $S$의 단위벡터 $x=\sum_{i\le k}c_iv_i$에서 $x^{\mathsf H}Hx=\sum_{i\le k}\lambda_i|c_i|^2\ge\lambda_k$이므로 오른쪽은 $\lambda_k$ 이상이다. 거꾸로 임의의 $k$차원 $S$와 $T=\operatorname{span}(v_k,\dots,v_n)$은 차원의 합이 $n+1$이므로 단위벡터 $x$를 공유하고, 그 $x$에서 $x^{\mathsf H}Hx=\sum_{i\ge k}\lambda_i|c_i|^2\le\lambda_k$이다. 따라서 어떤 $S$에서도 최솟값은 $\lambda_k$ 이하이다.
\end{proof}

\begin{corollary}[Weyl]
\label{cor:weyl}
에르미트 행렬 $H$와 $E$에 대해 $|\lambda_k(H+E)-\lambda_k(H)|\le\|E\|_2$이다. 또 모든 $m\times n$ 행렬 $A$, $E$에 대해 $|\sigma_k(A+E)-\sigma_k(A)|\le\|E\|_2$이다.
\end{corollary}

\begin{proof}
모든 단위벡터 $x$에서 $|x^{\mathsf H}Ex|\le\|E\|_2$이므로 $x^{\mathsf H}(H+E)x\le x^{\mathsf H}Hx+\|E\|_2$이고 \cref{thm:courant-fischer}의 최대최소를 양변에 취하면 $\lambda_k(H+E)\le\lambda_k(H)+\|E\|_2$이다. $H$와 $H+E$의 역할을 바꾸면 반대쪽 부등식을 얻는다. 특이값에 대해서는 $\sigma_k(A)=\max_{\dim S=k}\min_{x\in S,\|x\|=1}\|Ax\|_2$라는 같은 꼴의 표현이 \cref{thm:courant-fischer}를 $A^{\mathsf H}A$에 쓰면 나오고 $\|(A+E)x\|_2\le\|Ax\|_2+\|E\|_2$에서 같은 논증이 성립한다.
\end{proof}

이 결과는 1912년 Weyl이 미분방정식의 고유값 분포를 연구하며 얻은 부등식에서 이름을 따왔다\cite{weyl1912}. \cref{cor:weyl}을 계산의 관점에서 읽으면 특이값의 \textbf{절대} 오차는 행렬의 섭동 크기를 넘지 않는다. 역방향으로 안정한 알고리즘은 $\|E\|_2\approx u\|A\|_2$인 $A+E$의 특이값을 정확히 계산하므로 특이값의 절대 오차는 $u\sigma_1$ 정도이다. \cref{fig:svd-err}에서 본 것이 바로 이것이고 \cref{alg:rank-pinv}의 허용오차를 $\sigma_1$에 비례하게 잡은 것도 그래서이다.

\subsection{Non-Hermitian Matrices and the Bauer--Fike Theorem}

에르미트가 아닌 행렬의 고유값은 훨씬 민감할 수 있다. 크기가 $n$인 Jordan 블록 $J_n$, 곧 대각선 바로 위에만 $1$이 있는 행렬의 고유값은 모두 $0$이지만 왼쪽 아래 구석에 $\epsilon$을 더하면 특성다항식이 $\lambda^n-\epsilon$이 되어 고유값이 반지름 $\epsilon^{1/n}$인 원 위의 $n$개 점으로 뛴다. \cref{fig:jordan}은 $n=16$, $\epsilon=10^{-16}$에서 stardust의 \code{eigvals}가 계산한 고유값이다. $10^{-16}$의 섭동이 고유값을 $0.1$만큼 옮겼고 계산한 크기는 $0.1$과 $1.1\times10^{-14}$ 안쪽에서 맞는다. 섭동 전의 $J_{16}$은 이미 위삼각이어서 고유값이 정확히 $0$으로 나온다. 대각화 가능한 행렬에서는 이 민감도를 고유벡터 행렬의 조건수로 잴 수 있다.

\begin{figure}[htbp]
\centering
\begin{tikzpicture}
\begin{axis}[sdplot,axis equal image,xlabel={$\operatorname{Re}\lambda$},ylabel={$\operatorname{Im}\lambda$},xmin=-0.13,xmax=0.13,ymin=-0.13,ymax=0.13,width=0.45\textwidth,height=0.45\textwidth,xtick={-0.1,0,0.1},ytick={-0.1,0,0.1}]
\addplot[black!40,domain=0:360,samples=90,forget plot] ({0.1*cos(x)},{0.1*sin(x)});
\addplot[only marks,mark=*,mark size=1.6pt,black] table[x=re,y=im]{figures/04/jordan.dat};
\addlegendentry{$J_{16}+10^{-16}e_{16}e_1^{\mathsf T}$}
\addplot[only marks,mark=square,mark size=2.5pt,black!60] coordinates {(0,0)};
\addlegendentry{$J_{16}$}
\end{axis}
\end{tikzpicture}
\caption{Jordan 블록 $J_{16}$과 구석에 $10^{-16}$을 더한 행렬의 고유값(stardust \code{eigvals}). 회색 원은 반지름 $10^{-1}=(10^{-16})^{1/16}$이다.}
\label{fig:jordan}
\end{figure}

\begin{theorem}[Bauer--Fike]
\label{thm:bauer-fike}
$A=X\Lambda X^{-1}$이 대각화 가능하고 $\mu$가 $A+E$의 고유값이면 $\min_i|\mu-\lambda_i|\le\kappa_p(X)\|E\|_p$이다.
\end{theorem}

\begin{proof}
$\mu$가 $A$의 고유값이면 자명하므로 그렇지 않다고 하자. $A+E-\mu I$가 특이이므로 $X^{-1}(A+E-\mu I)X=(\Lambda-\mu I)+X^{-1}EX$도 특이이고 $\Lambda-\mu I$가 가역이므로 $I+(\Lambda-\mu I)^{-1}X^{-1}EX$가 특이이다. 그러면 $\|(\Lambda-\mu I)^{-1}X^{-1}EX\|_p\ge1$이어야 한다. 노름이 $1$보다 작은 행렬 $M$에 대해 $I+M$은 가역이기 때문이다($(I+M)x=0$이면 $\|x\|=\|Mx\|<\|x\|$). 대각행렬 $(\Lambda-\mu I)^{-1}$의 $p$-노름은 $1/\min_i|\lambda_i-\mu|$이므로 $1\le\|X^{-1}\|_p\|E\|_p\|X\|_p/\min_i|\lambda_i-\mu|$이고 이를 정리하면 주장이다.
\end{proof}

에르미트 행렬은 유니터리 행렬로 대각화되어 $\kappa_2(X)=1$이므로 \cref{thm:bauer-fike}는 \cref{cor:weyl}과 같은 크기의 한계를 준다\cite{bauer1960fike}. 고유벡터들이 거의 평행해 $\kappa(X)$가 크면 고유값은 민감하고 Jordan 블록처럼 대각화가 안 되는 극한에서는 섭동의 $1/n$제곱으로 움직인다. 대칭 고유값 알고리즘(54장)과 비대칭 고유값 알고리즘(55장)이 보장하는 정확도가 크게 다른 것도 이 때문이다.

\subsection{Perturbation of Subspaces and Least Squares}

고유값과 특이값에 대응하는 부분공간, 그리고 최소제곱 해가 섭동에 어떻게 반응하는지도 알아 두어야 한다. 이 결과들은 Wedin이 체계적으로 정리했다\cite{wedin1972}. 그 가운데 계산과 가장 직접 닿는 것은 최소제곱 해의 섭동이다. 열 랭크가 꽉 찬 $A$에 대한 문제 $\min\|Ax-b\|_2$에서 $A$와 $b$를 상대적으로 $\epsilon$만큼 바꾸면 해의 상대 변화는 일차까지
\[
\epsilon\Bigl(\kappa_2(A)+\kappa_2(A)^2\frac{\|r\|_2}{\|A\|_2\|x\|_2}\Bigr)
\]
의 크기로 상수 배 안에서 제한된다. 여기서 $r=b-Ax$는 잔차이다. 증명의 줄거리는 이렇다. 정규방정식 $A^{\mathsf H}(b-Ax)=0$을 섭동해 미분하면 $\Delta x=A^+(\Delta b-\Delta Ax)+(A^{\mathsf H}A)^{-1}\Delta A^{\mathsf H}r$이고 첫 항이 $\|A^+\|_2=1/\sigma_n$을 통해 $\kappa_2$를, 둘째 항이 $\|(A^{\mathsf H}A)^{-1}\|_2=1/\sigma_n^2$을 통해 $\kappa_2^2\|r\|$을 낸다. 잔차가 크면 조건수의 \textbf{제곱}이 나타난다. 그런데 정규방정식을 직접 풀면 잔차와 상관없이 $\kappa_2(A)^2$이 나타나므로 문제 자체보다 나쁜 알고리즘이 된다. 정확한 정리와 2차 항까지의 증명은 52장에서 직교 분해와 함께 다룬다. 섭동 이론 전반은 \cite{stewart1990perturbation}이 표준적인 참고서이다.

\section{The Schur Decomposition}
\label{sec:schur}

\subsection{The Theorem}

모든 정사각행렬이 대각화되지는 않지만 유니터리 닮음으로 \textbf{삼각화}할 수는 있다. 수치 계산에서 고유값을 다룰 때 가장 안정한 출발점이 이 삼각화이다.

\begin{theorem}[Schur]
\label{thm:schur}
모든 $n\times n$ 복소 행렬 $A$에 대해 유니터리 $Q$와 위삼각 $T$가 있어 $A=QTQ^{\mathsf H}$이다. $T$의 대각 성분은 $A$의 고유값이며 순서는 원하는 대로 정할 수 있다.
\end{theorem}

\begin{proof}
$n$에 대한 귀납법을 쓴다. 대수학의 기본정리에서 $A$는 고유값 $\lambda$와 단위 고유벡터 $q_1$을 가진다. $q_1$을 유니터리 $Q_1=[q_1\ \widetilde Q]$로 넓히면 $Q_1^{\mathsf H}AQ_1$의 첫 열은 $(\lambda,0,\dots,0)^{\mathsf T}$이므로
\[
Q_1^{\mathsf H}AQ_1=\begin{bmatrix}\lambda&*\\0&A_2\end{bmatrix}
\]
이다. $(n-1)\times(n-1)$ 행렬 $A_2$에 귀납 가정을 써서 $A_2=Q_2T_2Q_2^{\mathsf H}$를 얻고 $Q=Q_1\operatorname{diag}(1,Q_2)$로 두면 $Q^{\mathsf H}AQ$는 위삼각이다. 대각 성분이 고유값인 것은 닮음이 특성다항식을 보존하고 삼각행렬의 특성다항식이 $\prod(t_{ii}-\lambda)$이기 때문이다. 각 단계에서 어느 고유값을 먼저 잡을지 고를 수 있으므로 순서도 자유롭다.
\end{proof}

$A$가 정규행렬($AA^{\mathsf H}=A^{\mathsf H}A$)이면 $T$도 정규인 위삼각 행렬이고 그런 행렬은 대각행렬뿐이다. $T^{\mathsf H}T$와 $TT^{\mathsf H}$의 $(1,1)$ 원소를 견주면 $|t_{11}|^2=\sum_j|t_{1j}|^2$이므로 첫 행의 비대각 원소가 $0$이고 남은 블록에 같은 논증을 되풀이하면 된다. 따라서 정규행렬은 유니터리로 대각화된다. 실수 행렬은 켤레 복소 고유값 쌍 때문에 실수 범위에서 삼각화되지 않지만 대각선에 $1\times1$과 $2\times2$ 블록이 놓인 \term{실수 Schur 형식}{real Schur form}을 직교 행렬로 얻을 수 있다. 증명은 \cref{thm:schur}와 같은 귀납인데 복소 고유값 $\lambda=\alpha+i\beta$($\beta\ne0$)를 만나면 고유벡터의 실수부와 허수부가 생성하는 $2$차원 실수 불변 부분공간을 한 번에 떼어 내 $2\times2$ 블록을 만든다. 고유값 알고리즘은 바로 이 분해를 반복법으로 계산하며 그 과정은 55장에서 본다.

\subsection{The Sylvester Equation and the Bartels--Stewart Algorithm}

Schur 분해를 쓰는 대표적인 예가 Sylvester 방정식이라 부르는 행렬 방정식 $AX+XB=C$의 풀이이다($A$는 $m\times m$, $B$는 $n\times n$). $A=URU^{\mathsf H}$와 $B=WSW^{\mathsf H}$를 Schur 분해하고 $Y=U^{\mathsf H}XW$, $F=U^{\mathsf H}CW$로 바꾸면 방정식은 $RY+YS=F$가 된다. $R$과 $S$가 위삼각이므로 $Y$의 $k$번째 열 $y_k$를 보면
\[
(R+s_{kk}I)\,y_k=f_k-\sum_{j<k}s_{jk}\,y_j
\]
이고 오른쪽은 앞쪽 열들로 정해지므로 열을 왼쪽부터 차례로 삼각 방정식 하나씩 풀면 된다. 이것이 Bartels--Stewart 알고리즘이다\cite{bartels1972stewart}(\cref{alg:bartels-stewart}).

\begin{proposition}
\label{prop:sylvester-unique}
$AX+XB=C$가 모든 $C$에 대해 유일한 해를 가질 필요충분조건은 $A$의 고유값 $\lambda_i$와 $B$의 고유값 $\mu_j$가 모든 쌍에서 $\lambda_i+\mu_j\ne0$인 것이다.
\end{proposition}

\begin{proof}
위의 변환은 가역이므로 $RY+YS=F$를 보면 된다. 열마다의 삼각 방정식의 계수행렬 $R+s_{kk}I$의 대각 성분은 $r_{ii}+s_{kk}=\lambda_i+\mu_k$이다. 모두 $0$이 아니면 모든 열이 차례로 유일하게 풀린다. 어떤 쌍에서 $0$이면 그 연산자가 특이임을 \cref{sec:kronecker}의 Kronecker 꼴(\cref{prop:kron})이 보여 준다. 고유값이 $\lambda_i+\mu_j$이므로 그중 하나가 $0$이면 유일한 해가 없다.
\end{proof}

\begin{algorithm}[htbp]
\begin{algorithmic}[1]
\Require $A$ ($m\times m$), $B$ ($n\times n$), $C$ ($m\times n$)
\Ensure $AX+XB=C$의 해 $X$
\State $A=URU^{\mathsf H}$, $B=WSW^{\mathsf H}$ (Schur, \code{gees}) \Comment{서로 독립: 큰 행렬이면 동시에}
\State $F\gets U^{\mathsf H}CW$ \Comment{\code{gemm} 두 번}
\For{$k=1,\dots,n$} \Comment{실제로는 블록 단위(\code{trsyl3})}
  \State $g\gets\alpha f_k-\sum_{j<k}s_{jk}y_j$ \Comment{$\alpha$는 처음에 $1$}
  \State $(R+s_{kk}I)\,y_k=\beta g$를 뒤로 대입해 푼다 ($\beta\le1$) \Comment{넘침을 막는 척도}
  \State \textbf{if} $\beta<1$ \textbf{then} $y_1,\dots,y_{k-1}$에 $\beta$를 곱하고 $\alpha\gets\alpha\beta$
\EndFor
\State \textbf{if} 어떤 $|r_{ii}+s_{kk}|$가 너무 작다 \textbf{then} 오류 \Comment{고유값이 겹치거나 가깝다}
\State \Return $X=U(Y/\alpha)W^{\mathsf H}$ \Comment{\code{gemm} 두 번}
\end{algorithmic}
\caption{Bartels--Stewart 알고리즘 (stardust \code{solve\_sylvester})}
\label{alg:bartels-stewart}
\end{algorithm}

비용은 Schur 분해 두 번이 $O(m^3+n^3)$, 변환이 $O(mn(m+n))$, 삼각 풀이가 $O(mn(m+n))$이다. 방정식을 미지수 $mn$개짜리 보통의 연립방정식으로 바꾸어 풀면 $O((mn)^3)$이 드니 차이가 크다. stardust는 \cref{alg:bartels-stewart}를 LAPACK으로 짠다(\code{csrc/linalg/ops/special.cpp}). 두 Schur 분해는 서로 독립이므로 행렬이 크면 두 스레드에서 동시에 돌리고 같은 크기의 작업 공간을 쓰므로 결과 비트는 순서대로 돌린 것과 같다. 삼각 풀이는 블록 단위의 수준 3 루틴 \code{trsyl3}(LAPACK 3.10)를 쓰는데 해의 원소가 넘치지 않도록 우변 전체에 곱할 척도 $\alpha\le1$을 함께 돌려주므로 마지막에 $\alpha$로 나눈다. \code{trsyl3}가 $A$와 $-B$의 고유값이 겹치거나 너무 가까워 섭동한 값을 썼다고 알리면 stardust는 답을 돌려주지 않고 오류를 낸다. \cref{prop:sylvester-unique}에서 보듯 그런 방정식은 해가 없거나 조건수가 매우 크다. 실수 행렬은 실수 Schur 형식의 $2\times2$ 블록을 그대로 다룬다.

\cref{alg:bartels-stewart}의 4--6행은 척도를 누적한다. 어떤 열을 풀다 넘칠 것 같으면 그 열의 우변에 $\beta<1$을 곱한다. 그러면 앞서 푼 열들도 같은 척도로 맞춰야 $RY+YS=\alpha F$가 유지되므로 그 열들에 $\beta$를 곱하고 전체 척도를 $\alpha\beta$로 줄인다. 실수 Schur 형식에서 $S$에 $2\times2$ 블록이 있으면 그 두 열은 서로 엮이므로 $2m$개의 미지수를 한 번에 푼다. 해가 유일해도 문제의 조건이 좋다는 보장은 없다. Sylvester 연산자 $X\mapsto AX+XB$의 크기를 재는 양은
\[
\operatorname{sep}(A,-B)=\min_{X\ne0}\frac{\|AX+XB\|_F}{\|X\|_F}=\sigma_{\min}(I\otimes A+B^{\mathsf T}\otimes I)
\]
이고 해의 상대 섭동은 대략 $\|C\|_F/(\operatorname{sep}\cdot\|X\|_F)$배로 커진다($C$만 섭동할 때이고 $A$와 $B$도 섭동하면 $(\|A\|_F+\|B\|_F)/\operatorname{sep}$에 비례하는 항이 더해진다). $A$와 $B$가 정규이면 $\operatorname{sep}=\min_{i,j}|\lambda_i+\mu_j|$이지만 정규가 아니면 훨씬 작을 수 있다\cite{higham2002accuracy}. 또 Higham의 분석(그 책 16장)에 따르면 Bartels--Stewart는 잔차를 작게 만들지만 Sylvester 방정식의 역방향 오차는 잔차보다 클 수 있다. \cref{tab:sylvester}는 대각이 큰 무작위 행렬($A+mI$ 꼴)에서 \cref{alg:bartels-stewart}와 Kronecker 꼴의 $mn\times mn$ 연립방정식을 견준 것이다(\code{figures/04/measure.py}). 두 방법 모두 잔차가 반올림 수준이지만 시간은 $m=n=40$에서 이미 약 스물다섯 배 차이가 난다.

\begin{table}[htbp]
\begin{sdtable}{colspec={cllll}, row{1}={font=\small\bfseries}}
\toprule
$m=n$ & Bartels--Stewart & 잔차 & Kronecker 풀이 & 잔차 \\
\midrule
$20$ & $0.095$ms & $2.1\times10^{-15}$ & $0.40$ms & $1.2\times10^{-15}$ \\
$40$ & $0.31$ms & $2.9\times10^{-15}$ & $7.7$ms & $2.3\times10^{-15}$ \\
\bottomrule
\end{sdtable}
\caption{Sylvester 방정식의 두 풀이(stardust \code{solve\_sylvester}와 \code{kronsum} 꼴 행렬의 \code{solve}; 잔차는 $\|AX+XB-C\|_F/\|C\|_F$).}
\label{tab:sylvester}
\end{table}


\section{The Kronecker Product}
\label{sec:kronecker}

\subsection{Definition and Properties}

\begin{definition}[Kronecker 곱]
\label{def:kron}
$m\times n$ 행렬 $A$와 $p\times q$ 행렬 $B$의 \term{Kronecker 곱}{Kronecker product}은 $(i,j)$ 블록이 $a_{ij}B$인 $mp\times nq$ 블록 행렬 $A\otimes B$이다. 행렬 $X$의 열을 차례로 쌓은 벡터를 $\operatorname{vec}X$로 쓴다.
\end{definition}

원소로 쓰면 $(A\otimes B)_{ip+k,\,jq+l}=a_{ij}b_{kl}$이다($0$부터 센 색인, $0\le k<p$, $0\le l<q$). 행 색인 $ip+k$는 2장의 혼합 기수 $\phi(i,k)$(밑 $m,p$)이고 열 색인도 마찬가지이므로 \cref{thm:mixed-radix}에서 이 대응은 전단사이고 블록의 순서는 사전식이다.

\begin{proposition}
\label{prop:kron}
크기가 맞는 행렬들에 대해 다음이 성립한다.
\begin{enumerate}
\item $(A\otimes B)(C\otimes D)=(AC)\otimes(BD)$.
\item $\operatorname{vec}(AXB)=(B^{\mathsf T}\otimes A)\operatorname{vec}X$.
\item $A$가 $m\times m$이고 고유값이 $\lambda_i$, $B$가 $n\times n$이고 고유값이 $\mu_j$이면 $A\otimes B$의 고유값은 $\lambda_i\mu_j$이고 $I_n\otimes A+B\otimes I_m$의 고유값은 $\lambda_i+\mu_j$이다($i,j$의 모든 쌍).
\end{enumerate}
\end{proposition}

\begin{proof}
(i) 왼쪽의 $(i,k)$ 블록은 $\sum_ja_{ij}B\,c_{jk}D=(\sum_ja_{ij}c_{jk})BD=(AC)_{ik}BD$이다. (ii) $AXB$의 $k$번째 열은 $AXb_k=\sum_jb_{jk}Ax_j$이고($x_j$는 $X$의 $j$번째 열) 이는 $B^{\mathsf T}\otimes A$의 $k$번째 블록 행 $[b_{1k}A\ \cdots\ b_{nk}A]$를 $\operatorname{vec}X$에 곱한 것과 같다. (iii) Schur 분해 $A=Q_AT_AQ_A^{\mathsf H}$, $B=Q_BT_BQ_B^{\mathsf H}$와 (i)에서
\[
A\otimes B=(Q_A\otimes Q_B)(T_A\otimes T_B)(Q_A\otimes Q_B)^{\mathsf H}
\]
이고 $Q_A\otimes Q_B$는 (i)에서 유니터리이며 $T_A\otimes T_B$는 위삼각이고 대각 성분이 $\lambda_i\mu_j$이다. 합의 경우는 $I\otimes A+B\otimes I=(Q_B\otimes Q_A)(I\otimes T_A+T_B\otimes I)(Q_B\otimes Q_A)^{\mathsf H}$이고 가운데 행렬이 위삼각이며 대각 성분이 $\lambda_i+\mu_j$라는 데서 같다.
\end{proof}

\cref{prop:kron}의 (ii)를 쓰면 Sylvester 방정식 $AX+XB=C$는 $(I\otimes A+B^{\mathsf T}\otimes I)\operatorname{vec}X=\operatorname{vec}C$라는 보통의 연립방정식이 되고 (iii)에서 그 계수행렬의 고유값은 $\lambda_i(A)+\mu_j(B)$이다. 이것이 \cref{prop:sylvester-unique}의 나머지 절반이다. 다만 \cref{tab:sylvester}에서 본 대로 $mn\times mn$ 연립방정식을 그대로 풀면 비용이 $O((mn)^3)$이므로 실제 계산은 Bartels--Stewart처럼 구조를 살린다.

\begin{proposition}
\label{prop:kron-more}
$A$가 $m\times m$, $B$가 $n\times n$이면 $\det(A\otimes B)=\det(A)^n\det(B)^m$이고 둘 다 가역이면 $(A\otimes B)^{-1}=A^{-1}\otimes B^{-1}$이다. $A=L_AL_A^{\mathsf H}$, $B=L_BL_B^{\mathsf H}$가 Cholesky 분해이면 $A\otimes B=(L_A\otimes L_B)(L_A\otimes L_B)^{\mathsf H}$이고 $L_A\otimes L_B$는 대각이 양수인 하삼각이다. 특이값은 $\sigma_i(A)\sigma_j(B)$이다.
\end{proposition}

\begin{proof}
모두 \cref{prop:kron}의 (i)에서 나온다. $A\otimes B=(A\otimes I_n)(I_m\otimes B)$이고 $I_m\otimes B$는 $B$를 대각에 $m$번 놓은 블록 대각이므로 행렬식이 $\det(B)^m$, $A\otimes I_n$은 행과 열의 순서를 같은 순열로 바꾸면 $I_n\otimes A$가 되어 행렬식이 $\det(A)^n$이다. 역행렬은 $(A\otimes B)(A^{-1}\otimes B^{-1})=I\otimes I$에서, Cholesky는 $(L_A\otimes L_B)(L_A^{\mathsf H}\otimes L_B^{\mathsf H})$와 $(X\otimes Y)^{\mathsf H}=X^{\mathsf H}\otimes Y^{\mathsf H}$에서, 하삼각임은 블록 $(i,j)$가 $l_{ij}L_B$이고 $i<j$이면 $0$이라는 데서 나온다. 특이값은 두 특이값 분해의 Kronecker 곱 $(U_A\otimes U_B)(\Sigma_A\otimes\Sigma_B)(V_A\otimes V_B)^{\mathsf H}$을 보면 된다.
\end{proof}

행렬 $I_n\otimes A+B\otimes I_m$을 $A$와 $B$의 \term{Kronecker 합}{Kronecker sum}이라 하며 격자 위의 라플라시안이 그 대표적인 예이다. 길이 $m$인 선 위의 이산 라플라시안 $L_m$(대각이 $2$, 바로 옆이 $-1$인 삼중대각 행렬)으로 $m\times n$ 격자의 라플라시안을 만들면 $I_n\otimes L_m+L_n\otimes I_m$이 되고 \cref{prop:kron}에서 그 고유값은 두 선의 고유값의 합이다. 격자 위의 편미분방정식과 영상 처리에서 쓰는 빠른 풀이법 상당수가 이 분리 구조를 이용한다.

\subsection{Computing Kronecker Products and Sums}

stardust의 \code{kron}과 \code{kronsum}은 결과를 정의에서 바로 쓴다(\code{csrc/linalg/ops/basic.cpp}). \cref{alg:kron}은 열 우선 결과에 두 피연산자의 원소 곱을 놓는다.

\begin{algorithm}[htbp]
\begin{algorithmic}[1]
\Require $A$ ($m\times n$), $B$ ($p\times q$) (배치 축은 브로드캐스트, 1차원 $(N,)$은 행 $(1,N)$)
\Ensure $C=A\otimes B$ ($mp\times nq$, 열 우선)
\State $mp$, $nq$를 넘침 검사와 함께 계산 \Comment{넘치면 ``array is too big''}
\State $A$, $B$를 열 우선 작업 버퍼로 옮긴다 (dtype 변환 포함) \Comment{한 번만}
\For{$j=0,\dots,n-1$; $l=0,\dots,q-1$} \Comment{결과의 열 $jq+l$}
  \For{$i=0,\dots,m-1$}
    \State $C[ip:ip+p,\ jq+l]\gets a_{ij}\cdot B[:,l]$ \Comment{$a$를 왼쪽에: 복소 곱의 비트}
  \EndFor
\EndFor
\end{algorithmic}
\caption{Kronecker 곱 (stardust \code{kron})}
\label{alg:kron}
\end{algorithm}

\cref{alg:kron}의 안쪽 루프는 $B$의 한 열에 스칼라 하나를 곱해 결과의 연속한 $p$개 자리에 쓰므로 읽기와 쓰기가 모두 연속이다. 원소 곱 $mnpq$번과 같은 수의 쓰기가 드는데 결과가 크므로 계산보다 메모리 쓰기가 비용을 정한다. 1행의 넘침 검사는 2장의 \cref{prop:mul-overflow}이다. 브로드캐스트된 뷰에서는 $mp$나 $nq$가 64비트를 넘을 수 있기 때문이다. 5행에서 $a_{ij}$를 왼쪽에 두는 것은 numpy의 \code{kron}과 같은 비트를 내기 위해서이다. 실수 곱은 교환법칙이 비트까지 성립하지만 복소 곱은 실수부와 허수부를 만드는 연산의 순서가 달라 끝 비트가 다를 수 있다. 배치 축은 stardust의 다른 행렬 함수들처럼 브로드캐스트한다. numpy의 \code{kron}은 3차원 이상에서 모든 축에 Kronecker 곱을 하므로 그 경우에는 결과가 다르다.

Kronecker 합은 곱으로 계산하지 않는다. $I_n\otimes A+B\otimes I_m$의 원소는 $a$나 $b$나 $a+b$ 가운데 하나이므로 결과에 바로 놓는다(\cref{fig:kronsum}). 열 $J=jm+q'$에는 $A$의 $q'$번째 열이 블록 행 $j$에 들어간다. 또 블록 행 $r$마다 자리 $rm+q'$에 $b_{rj}$가 들어가는데 $r=j$이면 그 자리에서 $a_{q'q'}+b_{jj}$를 만든다(\cref{alg:kronsum}).

\begin{figure}[htbp]
\centering
\begin{tikzpicture}[x=0.62cm,y=0.62cm,font=\scriptsize]
\foreach \bi in {0,1}{\foreach \bj in {0,1}{
  \draw[thick] (3*\bj,-3*\bi) rectangle (3*\bj+3,-3*\bi-3);
}}
\foreach \b in {0,1}{
  \fill[gray!25] (3*\b,-3*\b) rectangle (3*\b+3,-3*\b-3);
  \node at (3*\b+2.25,-3*\b-0.75) {\large $A$};
}
\foreach \bi in {0,1}{\foreach \bj in {0,1}{
  \foreach \k in {0,1,2}{
    \node[circle,draw,inner sep=1pt,fill=white] at (3*\bj+\k+0.5,-3*\bi-\k-0.5) {\tiny $b$};
  }
  \node[anchor=south] at (3*\bj+1.5,-3*\bi+0.02) {};
}}
\node[anchor=west] at (6.4,-0.8) {회색 블록: $A$ (블록 대각)};
\node[anchor=west] at (6.4,-1.7) {원: $b_{rj}$ (블록 $(r,j)$의 대각)};
\node[anchor=west] at (6.4,-2.6) {회색 위의 원: $a_{qq}+b_{jj}$};
\node[anchor=west] at (6.4,-3.5) {나머지: $+0$};
\end{tikzpicture}
\caption{$m=3$, $n=2$일 때 $I_2\otimes A+B\otimes I_3$의 구조. 모든 원소가 $a$, $b$, $a+b$, $0$ 가운데 하나이다.}
\label{fig:kronsum}
\end{figure}

\begin{algorithm}[htbp]
\begin{algorithmic}[1]
\Require $A$ ($m\times m$), $B$ ($n\times n$)
\Ensure $K=I_n\otimes A+B\otimes I_m$ ($mn\times mn$, 열 우선)
\For{$J=0,\dots,mn-1$} \Comment{큰 결과는 열 $64$개씩 병렬로}
  \State $(j,q)\gets(\lfloor J/m\rfloor,\ J\bmod m)$; 열 $J$를 $+0$으로 \Comment{크면 처음부터 \code{calloc}}
  \State $K[jm:jm+m,\ J]\gets A[:,q]$
  \For{$r=0,\dots,n-1$}
    \State \textbf{if} $r=j$ \textbf{then} $K[rm+q,J]\gets K[rm+q,J]+b_{rj}$ \textbf{else} $K[rm+q,J]\gets b_{rj}$
  \EndFor
\EndFor
\end{algorithmic}
\caption{Kronecker 합 (stardust \code{kronsum})}
\label{alg:kronsum}
\end{algorithm}

곱을 쓰면 특수값이 엉뚱하게 퍼진다. $I\otimes A$를 곱으로 만들면 $A$의 NaN이 항등행렬의 $0$과 곱해져 $\mathrm{NaN}\cdot0=\mathrm{NaN}$이 되고 모든 블록의 같은 자리로 퍼진다. $4\times4$ 행렬 $A$에 NaN 하나, $3\times3$ 행렬 $B$에 $\infty$ 하나를 넣어 보면 $\code{kron}(I,A)+\code{kron}(B,I)$는 NaN이 $20$개인데 \cref{alg:kronsum}의 결과는 수학적으로 맞는 $3$개이다(\code{figures/04/measure.py}). 또 덧셈이 대각의 $a_{qq}+b_{jj}$ 한 번뿐이므로 결과가 같은 이름의 희소 함수 \code{sparse.kronsum}과 비트까지 같다. 따로 빌드한 stardust에서 두 결과를 64비트 정수로 읽어 견주면 일치한다. 배치가 행렬 하나이고 결과가 $2^{18}$개 원소 이상이며 열이 $64$개보다 많으면 열 $64$개씩을 스레드에 나누어 쓴다. 복사만 하므로 스레드 수가 결과를 바꾸지 않는다.

\section{Positive Definite Matrices}
\label{sec:pd}

\subsection{Characterizations}

에르미트 행렬 $H$가 모든 $x\ne0$에서 $x^{\mathsf H}Hx>0$이면 \term{양정치}{positive definite}, $\ge0$이면 \term{양반정치}{positive semidefinite}라 한다. 공분산 행렬, 헤시안의 최솟값 조건, 그람 행렬 $A^{\mathsf H}A$, 핵 방법의 핵 행렬이 모두 이 부류에 속하므로 양정치성을 알아보는 여러 방법을 정리해 두자.

\begin{theorem}
\label{thm:pd}
에르미트 행렬 $H$에 대해 다음은 서로 동치이다.
\begin{enumerate}
\item $H$는 양정치이다.
\item $H$의 모든 고유값이 양수이다.
\item 대각 성분이 양수인 하삼각 $L$이 있어 $H=LL^{\mathsf H}$이다(Cholesky 분해).
\item $H$의 모든 선행 주소행렬식 $\det H_{1:k,1:k}$가 양수이다(Sylvester 판정법).
\end{enumerate}
\end{theorem}

\begin{proof}
(i)$\Rightarrow$(ii): $Hv=\lambda v$, $v\ne0$이면 $0<v^{\mathsf H}Hv=\lambda\|v\|_2^2$이다. (ii)$\Rightarrow$(i): 스펙트럼 정리로 $H=V\Lambda V^{\mathsf H}$이고 $x^{\mathsf H}Hx=\sum\lambda_i|(V^{\mathsf H}x)_i|^2>0$이다. (iii)$\Rightarrow$(i): $x^{\mathsf H}Hx=\|L^{\mathsf H}x\|_2^2$이고 $L$이 가역이므로 $x\ne0$이면 양수이다. (i)$\Rightarrow$(iv): 선행 주소행렬 $H_k$도 양정치이므로(그 부분 벡터에 $0$을 덧붙여 생각하면 된다) (ii)에서 고유값의 곱인 행렬식이 양수이다. (iv)$\Rightarrow$(iii): $n$에 대한 귀납법을 쓴다. $H$를
\[
H=\begin{bmatrix}H_{n-1}&h\\h^{\mathsf H}&\eta\end{bmatrix}
\]
로 나누면 귀납 가정에서 $H_{n-1}=L_{n-1}L_{n-1}^{\mathsf H}$이다. $l=L_{n-1}^{-1}h$와 $\gamma=\eta-l^{\mathsf H}l$로 두면
\[
H=\begin{bmatrix}L_{n-1}&0\\l^{\mathsf H}&1\end{bmatrix}\begin{bmatrix}I&0\\0&\gamma\end{bmatrix}\begin{bmatrix}L_{n-1}^{\mathsf H}&l\\0&1\end{bmatrix}
\]
이고 양변의 행렬식에서 $\det H=\gamma\det H_{n-1}$이다. 가정에서 두 행렬식이 양수이므로 $\gamma>0$이고 $\gamma=\sqrt\gamma\cdot\sqrt\gamma$로 쪼개 바깥 두 인수에 하나씩 붙이면 Cholesky 분해를 얻는다.
\end{proof}

증명의 마지막 단계에 나온 $\gamma=\eta-h^{\mathsf H}H_{n-1}^{-1}h$를 $H_{n-1}$에 대한 \term{Schur 보수}{Schur complement}라 한다. 블록 행렬
\[
\begin{bmatrix}P&Q\\Q^{\mathsf H}&R\end{bmatrix}
\]
이 양정치일 필요충분조건은 $P$와 Schur 보수 $R-Q^{\mathsf H}P^{-1}Q$가 모두 양정치인 것이다. $P$가 가역이면 블록 합동
\[
\begin{bmatrix}P&Q\\Q^{\mathsf H}&R\end{bmatrix}=\begin{bmatrix}I&0\\Q^{\mathsf H}P^{-1}&I\end{bmatrix}\begin{bmatrix}P&0\\0&R-Q^{\mathsf H}P^{-1}Q\end{bmatrix}\begin{bmatrix}I&P^{-1}Q\\0&I\end{bmatrix}
\]
이 성립하고 가역 행렬로 합동인 두 에르미트 행렬은 하나가 양정치이면 다른 하나도 양정치이므로(Sylvester의 관성 법칙의 쉬운 쪽) 주장이 나온다. 다변량 정규분포의 조건부 공분산이 바로 Schur 보수이므로 이 사실은 그 조건부 분포를 다룰 때 다시 나온다(X권).

\subsection{Testing Definiteness and Evaluating Quadratic Forms}

수치 계산에서는 (iii)이 가장 실용적인 판정법이다. 고유값을 구하는 것보다 훨씬 싸고 판정과 함께 이차 형식과 행렬식에 필요한 인수를 준다. \cref{thm:pd}의 증명이 그대로 알고리즘이 된다(\cref{alg:pd-cholesky}). 열을 하나씩 늘리며 Schur 보수 $\gamma$가 양수인지 본다.

\begin{algorithm}[htbp]
\begin{algorithmic}[1]
\Require 에르미트 $H$ ($d\times d$, 아래 삼각만 읽는다), 벡터 $x$
\Ensure 양정치가 아니면 ``실패'', 아니면 $Q=x^{\mathsf H}H^{-1}x$와 $\ln\det H$
\For{$j=1,\dots,d$}
  \State $\gamma\gets h_{jj}-\sum_{k<j}|l_{jk}|^2$ \Comment{Schur 보수}
  \State \textbf{if} $\gamma\le0$ \textbf{or} $\gamma$가 유한하지 않다 \textbf{then} \Return 실패
  \State $l_{jj}\gets\sqrt\gamma$; $\ell\ell\gets\ell\ell+2\ln l_{jj}$ \Comment{$\ln\det H=2\sum\ln l_{jj}$}
  \For{$i=j+1,\dots,d$} $l_{ij}\gets(h_{ij}-\sum_{k<j}l_{ik}\overline{l_{jk}})/l_{jj}$ (유한하지 않으면 실패)
  \EndFor
\EndFor
\State $Lz=x$를 앞으로 대입해 푼다; $Q\gets\|z\|_2^2$ \Comment{역행렬을 만들지 않는다}
\State \Return $Q$, $\ell\ell$
\end{algorithmic}
\caption{Cholesky로 하는 양정치 판정과 이차 형식}
\label{alg:pd-cholesky}
\end{algorithm}

\cref{alg:pd-cholesky}의 8행은 $x^{\mathsf H}H^{-1}x=x^{\mathsf H}L^{-\mathsf H}L^{-1}x=\|L^{-1}x\|_2^2$이라는 데서 나온다. 역행렬을 만들고 곱하는 것보다 연산이 적고 결과가 음수로 나올 수 없다. 4행의 $\ln\det H=2\sum_j\ln l_{jj}$는 행렬식을 직접 곱하면 넘치거나 아래로 넘치는 큰 차원에서도 안전하다. 비용은 Cholesky가 $d^3/6$번의 곱셈–덧셈($d^3/3$ flop), 대입이 $d^2/2$번이다. 블록화한 Cholesky 알고리즘과 그 안정성은 50장에서 다룬다. stardust의 \code{linalg.cholesky}는 3행의 실패를 오류로 알린다.

이차 형식의 정확도는 조건수에 걸려 있다. binary64로 분해하면 $L$은 $\|E\|\approx du\|H\|$인 $H+E$의 정확한 인수이므로 $Q$의 상대오차가 $\kappa(H)u$ 정도로 커질 수 있다. stardust의 다변량 정규분포는 이것을 피하려고 인수 $L$을 double-double로 만든다(\code{csrc/stats/dd\_linalg.hpp}, double-double 산술은 12장). 그러면 인수에서 오는 오차가 $\kappa(H)\cdot10^{-32}$ 정도로 내려간다. 다만 점마다 double-double 대입을 하면 느리므로 매개변수마다 한 번 $W=L^{-1}$을 double-double로 만들어 binary64로 반올림해 두고 점마다 $z=W(x-\mu)$를 평범한 내적으로 구하면서 그 오차의 엄밀한 한계를 함께 계산한다(\cref{alg:mvn-fast}). 한계가 예산 안에 들면 그 값을 쓰고 아니면 double-double 대입으로 다시 계산한다.

\begin{algorithm}[htbp]
\begin{algorithmic}[1]
\Require 매개변수마다: double-double 인수 $L$, $\widehat W=\mathrm{fl}(L^{-1})$과 행 노름; 점 $x$
\Ensure $Q=\|L^{-1}(x-\mu)\|_2^2$ (로그 밀도 기준 예산 $40u$ 안에서)
\State $r\gets x-\mu$; $z\gets\widehat Wr$ \Comment{행마다 \code{fma} 사슬 $8$개, 정해진 트리}
\State $e_i\gets\gamma_{m_i}\bigl(|\widehat w_{ii}||r_i|+\|\widehat w_{i,<i}\|_2\|r_{<i}\|_2\bigr)+\cdots$ \Comment{Cauchy--Schwarz로 한 번에}
\State $Q\gets\|z\|_2^2$, $\varepsilon_Q\gets\sum_i(2|z_i|+e_i)e_i+\gamma Q$
\State \textbf{if} $\varepsilon_Q\le40u\cdot\max(1,|\text{로그 밀도}|)$ \textbf{then} \Return $Q$
\State \Return double-double 전진 대입으로 계산한 $Q$ \Comment{조건이 나쁘거나 상쇄가 큰 점}
\end{algorithmic}
\caption{다변량 정규 이차 형식의 빠른 경로 (stardust \code{fast\_solve}, \code{fast\_norm2})}
\label{alg:mvn-fast}
\end{algorithm}

\cref{alg:mvn-fast}의 한계는 내적의 표준적인 오차 한계로 얻는다. 행 $i$의 내적은 갈래 $8$개의 \code{fma} 사슬을 정해진 트리로 합치므로 어떤 항이든 반올림을 $m_i$번 이하 거치고 그 오차는 $\gamma_{m_i}\sum_k|\widehat w_{ik}||r_k|$ 이하이다($\gamma_m=mu/(1-mu)$, 11장). 합을 다시 읽지 않으려고 대각 밖의 부분은 Cauchy--Schwarz로 $\|\widehat w_{i,<i}\|_2\|r_{<i}\|_2$ 이하로 누르고 행 노름은 매개변수마다 미리 구해 둔다. 여기에 $\widehat W$ 자신의 반올림과 $W$를 만든 double-double 대입의 오차가 더해진다(2행의 생략한 항). 예산 $40u$는 계약 $10^{-14}$ 안에서 정한 값이고 조건이 나쁜 공분산에서는 한계가 커져 대부분의 점이 5행으로 간다. 양정치가 아니면 로그 밀도는 NaN이고 오류는 내지 않는다. \cref{fig:mvn}은 Hilbert 행렬을 공분산으로 둔 $d$차원 정규분포의 로그 밀도를 무작위 $x$에서 계산한 것이다. 참값은 binary64 Hilbert 행렬의 이차 형식과 행렬식을 정확한 유리수로 구해 만들었다(\code{figures/04/mvn\_logpdf.py}). binary64 Cholesky의 상대오차는 $\kappa u$를 따라 $d=11$($\kappa=5\times10^{14}$)에서 $10^{-3}$에 이른다. 반면 따로 빌드한 stardust의 \code{MultivariateNormal.logpdf}는 모든 $d$에서 $2\times10^{-16}$ 아래이다. 조건이 나쁜 오른쪽 점들에서는 \cref{alg:mvn-fast}의 한계가 예산을 넘어 double-double 대입이 답을 낸다.

같은 도구가 행렬 분포에서도 쓰인다. Wishart 분포의 밀도는 점 $X$마다 $X$의 double-double Cholesky를 시도하고 이 시도가 곧 지지집합의 검사가 된다. 실패하면 $X$가 양정치가 아니어서 지지집합 밖이므로 밀도가 $0$이다. 역 Wishart의 $\operatorname{tr}(\Psi X^{-1})$은 $X=L_XL_X^{\mathsf H}$, $\Psi=L_\Psi L_\Psi^{\mathsf H}$에서 $\operatorname{tr}(L_\Psi^{\mathsf H}L_X^{-\mathsf H}L_X^{-1}L_\Psi)=\|L_X^{-1}L_\Psi\|_F^2$이므로 역행렬 없이 삼각 풀이와 Frobenius 노름으로 계산한다(\code{csrc/bind/stats\_matrix.cpp}). 행렬 정규분포 $\mathcal{MN}(M,U,V)$의 공분산은 $V\otimes U$이고 다음 절의 \cref{prop:kron-more}에서 $\ln\det(V\otimes U)=p\ln\det U+n\ln\det V$, 이차 형식은 $\|L_U^{-1}(X-M)L_V^{-\mathsf H}\|_F^2$이 된다.

\begin{figure}[htbp]
\centering
\begin{tikzpicture}
\begin{loglogaxis}[sdplot,xlabel={$\kappa_2(H)$},ylabel={로그 밀도의 상대오차},xmin=5,xmax=1e16,ymin=1e-19,ymax=1e-1,ytick={1e-18,1e-14,1e-10,1e-6,1e-2},yticklabels={${\le}10^{-18}$,$10^{-14}$,$10^{-10}$,$10^{-6}$,$10^{-2}$}]
\addplot[black,mark=*,mark size=1.3pt] table[x=kappa,y=f64]{figures/04/mvn.dat};
\addlegendentry{binary64 Cholesky}
\addplot[black,dashed,mark=square*,mark size=1.3pt] table[x=kappa,y=lib]{figures/04/mvn.dat};
\addlegendentry{stardust}
\addplot[black!40,dotted,thick,domain=5:1e16,samples=2] {1.11e-16*x};
\addlegendentry{$\kappa u$}
\end{loglogaxis}
\end{tikzpicture}
\caption{Hilbert 행렬을 공분산으로 둔 $d=2,\dots,11$차원 정규분포의 로그 밀도 오차($10^{-18}$은 정확히 맞은 점).}
\label{fig:mvn}
\end{figure}

\section{Testing Matrix Code}
\label{sec:matrix-testing}

행렬 코드를 시험하는 기준은 세 가지이다. 하나는 정확한 유리수이다. 이 장의 측정은 binary64 입력을 정확한 유리수로 보고 해, 이차 형식, 행렬식, 노름을 정확히 계산해 참값으로 삼았다. 크기가 수십까지인 행렬에서는 이 방법이 충분히 빠르고 오차가 없다. 답을 미리 아는 행렬도 기준이 된다. 정수 벡터의 Householder 반사처럼 유리수 직교행렬로 특이값을 정해 둔 행렬(\cref{fig:svd-err}), 고유값을 손으로 아는 Jordan 블록(\cref{fig:jordan}), Hilbert 행렬처럼 조건수가 알려진 행렬이 그런 예이다. 마지막은 조건수에 맞춘 판정이다. \cref{thm:linsys-perturb}에서 보았듯 역방향으로 안정한 알고리즘도 $\kappa u$만큼 틀리므로 ``참값과 $u$ 안에서 같은가''가 아니라 ``잔차가 반올림 수준인가'', ``전방 오차가 $\kappa u$의 작은 배수인가''를 본다. stardust 자신의 시험도 같은 기준을 쓴다. 노름은 여러 레이아웃·dtype·축의 조합에서 느린 기준 구현과 견준다(\code{tests/linalg/test\_basic\_norms.py}). 그 밖에 스레드 수에 따라 비트가 같은지(\code{test\_parallel.py}), 희소 함수와 비트가 같은지(\code{test\_sparse\_counterparts.py}), 비정규수 특이값에서 \code{pinv}가 맞는지(\code{test\_tiny\_singular\_values.py})를 따로 본다. \cref{tab:matrix-tests}는 이 장의 측정을 모은 것이다. 모든 수치는 공개 저장소를 따로 받아 빌드한 stardust로 쟀다.

\begin{table}[htbp]
\begin{sdtable}{colspec={lX}, row{1}={font=\small\bfseries}}
\toprule
대상 & 기준과 결과 \\
\midrule
\code{vector\_norm} & 정확한 제곱합의 정수 제곱근; $2^{-1070}$--$2^{1010}$배 척도 모두 올바르게 반올림 \\
\code{matrix\_norm} & 시간(배치별), C·F 순서의 끝 비트 차이($40$개 중 $17$개, 설계가 받아들인 것) \\
\code{svdvals} & 정해진 특이값; 상대오차가 $u\sigma_1/\sigma_i$ 아래 \\
\code{solve} & Hilbert, 정확한 유리수 해; 오차 $\approx\kappa u$ \\
\code{eigvals} & Jordan 블록; 반지름 $0.1$의 원 위, $1.1\times10^{-14}$ 안 \\
\code{solve\_sylvester} & 잔차 $3\times10^{-15}$ 이하, Kronecker 풀이와 비교 \\
\code{kronsum} & \code{sparse.kronsum}과 비트 일치, NaN의 자리 \\
\code{MultivariateNormal.logpdf} & 정확한 이차 형식과 행렬식; $2\times10^{-16}$ 이하 \\
\bottomrule
\end{sdtable}
\caption{이 장의 측정과 시험(\code{figures/04/}의 스크립트).}
\label{tab:matrix-tests}
\end{table}

\section{APIs in This Chapter}
\label{sec:matrix-apis}

\cref{tab:matrix-apis}에 이 장의 커널을 쓰는 stardust의 API와 그 커널을 설명한 곳을 모았다. 분해 자체(LU 48장, Cholesky 50장, QR 51장, 고유값 54--55장, SVD 56장)와 조건수 어림(53장), 행렬 함수와 행렬 방정식(57장)은 V권에서 다루고, 47장은 이 장의 노름과 Kronecker 구현을 다시 쓰지 않고 가리킨다.

\begin{table}[htbp]
\begin{sdtable}{colspec={lXl}, row{1}={font=\small\bfseries}}
\toprule
API & 이 장에서 다룬 것 & 설명한 곳 \\
\midrule
\code{linalg.vector\_norm} & 두 단계 척도 맞춤, $p$-노름, 특수값 & \cref{alg:norm2} \\
\code{linalg.matrix\_norm} & 선 합의 방향과 띠, \code{'fro'}, 특이값 노름 & \cref{alg:line-norm} \\
\code{linalg.matrix\_rank}, \code{pinv} & 허용오차 $\max(m,n)\cdot2^{-52}\cdot\sigma_1$ & \cref{alg:rank-pinv} \\
\code{linalg.solve\_sylvester} & Bartels--Stewart, 척도, 가까운 고유값 & \cref{alg:bartels-stewart} \\
\code{linalg.kron}, \code{kronsum} & 정의에서 직접, 곱 없는 Kronecker 합 & \cref{alg:kron,alg:kronsum} \\
\code{stats.MultivariateNormal}, \code{Wishart}, \code{MatrixNormal} & 양정치 판정, 이차 형식, $\ln\det$, 빠른 경로의 오차 한계 & \cref{alg:pd-cholesky,alg:mvn-fast} \\
\code{linalg.cond} & $\sigma_1/\sigma_n$, 다른 \code{ord}의 노름 & \cref{sec:condition} \\
\bottomrule
\end{sdtable}
\caption{이 장의 커널을 쓰는 stardust API.}
\label{tab:matrix-apis}
\end{table}

\section{Concluding Remarks}

이 장은 행렬의 크기를 재는 노름에서 출발했지만 정리 대부분은 특이값 분해 위에 섰다. 노름의 동치성, 유도 노름의 공식, Eckart--Young, 조건수와 특이행렬까지의 거리, Weyl과 Bauer--Fike의 섭동 한계를 모두 특이값과 고유값으로 쓸 수 있다. 알고리즘은 이 정리들을 계산 순서로 옮긴 것이다. 넘치지 않는 노름은 척도의 동차성을 쓰고 수치적 랭크는 Weyl의 한계에 기댄다. Bartels--Stewart는 Schur 분해의 삼각 구조를 따라가며 양정치 판정은 Schur 보수의 귀납을 그대로 밟는다. 구현으로 내려오면 메모리 배치와 반올림이 끼어든다. 열 합을 띠로 묶을지, Kronecker 합을 곱으로 만들지 직접 놓을지, 이차 형식을 double-double로 계산할지에 따라 속도와 정확도, 특수값을 다루는 방식이 달라진다. 다음 장에서는 특수함수와 변환의 바탕이 되는 복소해석을 다룬다.

\chapterreferences
